%% IEEEtran skeleton for the manuscript. The body is produced MECHANICALLY from
%% manuscript.md via pandoc (see paper/ieee/README.md) rather than hand-converted,
%% so the math/tables/proofs carry over faithfully. This file is the correct,
%% compilable wrapper; it was NOT compiled in the authoring environment (no LaTeX
%% toolchain present) — run the build in paper/ieee/README.md to produce the PDF.
\documentclass[conference]{IEEEtran}
\IEEEoverridecommandlockouts

\usepackage{amsmath,amssymb,amsfonts}
\usepackage{amsthm}
\usepackage{algorithmic}
\usepackage{graphicx}
\usepackage{booktabs}
\usepackage{array}
\usepackage{longtable}
\usepackage{tabularx}
\usepackage{textcomp}
\usepackage{upquote}
\providecommand{\tightlist}{\setlength{\itemsep}{0pt}\setlength{\parskip}{0pt}}
\providecommand{\pandocbounded}[1]{#1}
\providecommand{\passthrough}[1]{#1}
\providecommand{\textquotesingle}{'}
% Shrink-to-fit for wide display equations so they never overflow a column.
\newcommand{\fiteq}[1]{\par\vspace{2pt}\noindent\centerline{%
  \resizebox{\ifdim\width>0.96\linewidth 0.96\linewidth\else \width\fi}{!}{%
  \ensuremath{\displaystyle #1}}}\vspace{2pt}\par}
\usepackage{listings}
\usepackage{xcolor}
\usepackage{url}
\usepackage[hidelinks]{hyperref}
\usepackage{cite}

\newtheorem{theorem}{Theorem}
\newtheorem{lemma}{Lemma}
\newtheorem{definition}{Definition}

% Compact code font so the Fig.1 ASCII dataflow fits a single column.
\lstset{basicstyle=\ttfamily\scriptsize,breaklines=true,breakatwhitespace=false,%
  columns=fullflexible,keepspaces=true,frame=single,xleftmargin=0pt}

\begin{document}

\title{Multi-Layered Defensive Architecture and Webhook Securitization for LLM
Deployment in Multi-Tenant Enterprise Systems}

\author{\IEEEauthorblockN{Mohammad Safari}
\IEEEauthorblockA{AIO Apex\\
Brighton, United Kingdom\\
Email: mo@aioapex.com}
\thanks{Code and data available at \url{https://github.com/mosafariuk/prompt-guard-2-frozen-head}}}

\maketitle

\begin{abstract}
Large Language Models exposed through webhook interfaces in multi-tenant enterprise
platforms inherit a hostile attack surface in which a single unauthenticated or
adversarial request can trigger prompt injection, cross-tenant data leakage, or
spoofed event delivery. This paper presents a dual-tier defensive architecture. A
multi-tenant V8-isolate edge (i)~authenticates every webhook via HMAC-SHA256 signatures
binding a tenant identifier into the signed message---providing EUF-CMA tenant-context
integrity with a multi-user bound independent of tenant count and key-rotation
multiplicity---and (ii)~screens payloads with deterministic single-pass Aho--Corasick
matching in a sub-millisecond CPU budget, \emph{failing open} to a self-hosted origin
deep-scanner for availability; interception events are (iii)~logged to a partitioned
PostgreSQL JSONB store on a non-blocking path. We formalize the threat model with STRIDE
extended by cross-boundary coupling and correct a widely-repeated misconception about
edge CPU limits. Our central empirical finding concerns the origin tier: Meta Prompt
Guard~2's 22.8\% out-of-distribution recall is a \emph{decision-head artifact}, not an
encoder limitation---a linear head on its frozen embeddings reaches \textbf{99.9\% OOD
recall at 0.7\% false-positive rate (AUC $\approx$ 0.999)} under a pre-registered,
leakage-controlled protocol on fresh disjoint data, and is deployed as the production
classifier. We treat the evaluation discipline as a first-class contribution:
pre-registered criteria, cosine-$\ge$0.95 cross-split deduplication, refusal of
unverifiable-license data, and two NULL verdicts held on 99.7\%+ recall before the
strict false-positive ceiling was met.
\end{abstract}

\begin{IEEEkeywords}
LLM security, prompt injection, multi-tenant isolation, edge computing, HMAC,
webhook authentication, PostgreSQL, JSONB, STRIDE.
\end{IEEEkeywords}

%% -----------------------------------------------------------------------------
%% BODY. Generate body.tex from manuscript.md (see README.md), then:
\section{I. Introduction}\label{i-introduction}

\subsection{I-A. Motivation}\label{i-a-motivation}

Enterprise SaaS platforms increasingly place Large Language Models
(LLMs) on the \emph{receiving} end of webhooks. A payment-orchestration
provider routes \passthrough{\lstinline!charge.disputed!} events into an
LLM that drafts a merchant response; a ticketing platform feeds inbound
\passthrough{\lstinline!issue.created!} payloads to an LLM triage agent;
a workflow-automation vendor lets tenants wire arbitrary third-party
events into model-backed actions. In each case the LLM is a
\emph{shared, multi-tenant} compute resource invoked by \emph{externally
originated, attacker-influenceable} HTTP requests. This composition is
qualitatively more dangerous than an interactive chatbot: the input is
machine-delivered rather than human-typed, it arrives at machine rates,
and a single event may carry both tenant A\textquotesingle s trusted
configuration and tenant B\textquotesingle s adversarial content into
the same model context.

Throughout this paper we use a \textbf{payment-gateway} as the running
example because it exhibits the full hazard profile --- high request
velocity, strong tenant-isolation requirements, regulatory exposure, and
a webhook interface that is, by construction, reachable by parties the
platform does not fully trust.

\subsection{I-B. Problem Statement}\label{i-b-problem-statement}

The prevailing defensive posture places LLM guardrails at the
\emph{application origin}: the request traverses the public network,
terminates at the platform\textquotesingle s backend, and only there is
it inspected --- for prompt injection, for signature validity, for
tenant scope. This arrangement has two structural defects.

\begin{enumerate}
\def\labelenumi{\arabic{enumi}.}
\tightlist
\item
  \textbf{Latency and cost are paid before rejection.} A malicious
  webhook consumes a full network round trip and origin compute before
  it is discarded. Under a flood, the origin absorbs the full load of
  traffic it will ultimately reject.
\item
  \textbf{The trust boundary and the inspection point coincide.} The
  component that must be protected is also the component performing the
  protection. A parser bug or an injection that survives the guardrail
  executes in the same trust domain as tenant data.
\end{enumerate}

A defense that inspects and authenticates requests \emph{before} they
reach the origin --- at the network edge, in a sandboxed runtime that
shares no state with tenant data --- addresses both. The engineering
question this paper answers is whether such an edge defense can be made
\textbf{cryptographically sound} (it must authenticate tenants, not
merely filter text), \textbf{computationally cheap enough} to run inside
a constrained edge runtime, and \textbf{observable} (every interception
durably logged) without any of these properties gating the request path.

\subsection{I-C. A Premise We Must Correct Up
Front}\label{i-c-a-premise-we-must-correct-up-front}

A recurring claim in edge-security folklore --- and in the original
specification of this work --- is that a "standard Cloudflare Worker
enforces a 50 ms CPU limit," and that fitting a firewall inside that
window is the central engineering challenge. \textbf{This is incorrect
as of 2026.} Our verification against primary Cloudflare documentation
establishes that the 50 ms figure is a \emph{legacy artifact}: it was
auto-applied to the deprecated "Bundled" usage model during the
migration to Standard pricing on 2024-03-01 {[}C-pricing{]}. The current
envelope (§III) is 10 ms of CPU per invocation on the Free plan, and on
paid plans a default of \textbf{30 seconds}, configurable to \textbf{5
minutes} {[}C-limits, C-changelog-2025{]}. Moreover, CPU time meters
\emph{active computation only}; time spent awaiting network, KV, or
database I/O is excluded {[}C-limits{]}.

Two consequences follow, and we treat them as first-class contributions
rather than footnotes. First, the design target for our deterministic
screening layer is \textbf{sub-millisecond CPU}, which sits three to
four orders of magnitude below even the Free-plan ceiling --- the
engineering story is \emph{headroom}, not survival. Second, the
"\textless50 ms" budget that genuinely matters is a \textbf{wall-clock,
end-to-end latency SLO}, a distinct quantity from the CPU-time limit.
Conflating the two --- as the folklore does --- produces both an
overstated challenge and an unfalsifiable evaluation. We keep them
separate throughout, and §III and §VII report them independently.

\subsection{I-D. Contributions}\label{i-d-contributions}

This paper makes the following contributions.

\begin{enumerate}
\def\labelenumi{\arabic{enumi}.}
\tightlist
\item
  \textbf{An edge-relocated, multi-layered defensive architecture} for
  LLM webhook ingestion in multi-tenant systems, in which
  authentication, screening, and logging execute in a V8-isolate edge
  runtime ahead of the origin (§III).
\item
  \textbf{A tenant-context-binding construction and its security
  argument.} We bind a tenant identifier into an HMAC-SHA256-signed
  message and prove, by reduction to HMAC\textquotesingle s existential
  unforgeability under chosen-message attack (EUF-CMA), that an attacker
  cannot present a validly-signed webhook attributed to a tenant it does
  not control (§IV).
\item
  \textbf{A deterministic screening layer with explicit complexity
  bounds.} Aho-Corasick multi-pattern matching in O(n+z) over all
  injection signatures in a single pass, plus O(n) structural and
  entropy scoring, with a worst-case CPU-budget accounting that
  establishes the sub-millisecond claim against the measured §III
  envelope (§V).
\item
  \textbf{A non-blocking threat-logging data path.} Interception events
  are written out-of-band to a partitioned PostgreSQL JSONB store after
  the response is returned, so that durable logging never appears on the
  request\textquotesingle s critical path (§VI).
\item
  \textbf{An empirical evaluation} reporting p50/p90/p99 latency (CPU
  and wall-clock reported separately) and detection efficacy ---
  false-positive and false-negative rates --- against a documented
  open-source prompt-injection corpus, together with the analytical
  latency budget that the measurements test (§VII).
\end{enumerate}

As a cross-cutting contribution, we \textbf{date-stamp and correct the
edge resource model} (I-C, §III), replacing a widely-repeated but
obsolete constraint with the current, primary-sourced envelope.

\subsection{I-E. Paper Organization}\label{i-e-paper-organization}

Section II formalizes the threat model with STRIDE and fixes the
prompt-injection and cross-tenant-leakage taxonomy. Section III
characterizes the V8-isolate execution envelope. Section IV develops the
cryptographic tenant-isolation construction and its proof. Section V
specifies the heuristic screening layer and its complexity. Section VI
details the asynchronous JSONB logging architecture. Section VII
presents the evaluation methodology and results. Sections VIII--X cover
related work, limitations, and conclusions.

\begin{center}\rule{0.5\linewidth}{0.5pt}\end{center}

\subsubsection{Citation keys used above (resolved in the References
section)}\label{citation-keys-used-above-resolved-in-the-references-section}

\begin{itemize}
\tightlist
\item
  \textbf{{[}C-pricing{]}} Cloudflare, "Workers Pricing,"
  developers.cloudflare.com/workers/platform/pricing (accessed
  2026-07-04).
\item
  \textbf{{[}C-limits{]}} Cloudflare, "Workers Limits,"
  developers.cloudflare.com/workers/platform/limits (accessed
  2026-07-04).
\item
  \textbf{{[}C-changelog-2025{]}} Cloudflare, "Higher CPU limits,"
  changelog 2025-03-25.
\end{itemize}

\begin{quote}
(\passthrough{\lstinline!paper/refs/verified-facts.md!}), Part 1 --- all
CONFIRMED at high confidence (3-0 votes).
\end{quote}

\section{II. Threat Model and
Taxonomy}\label{ii-threat-model-and-taxonomy}

We first fix the system and adversary model (§II-A), then decompose the
attack surface with STRIDE, applied with explicit cross-boundary
coupling to account for the code/data confusion intrinsic to LLMs
(§II-B). We then pin the prompt-injection taxonomy to the OWASP LLM Top
10 (§II-C) and formalize cross-tenant leakage (§II-D).

\subsection{II-A. System and Trust
Model}\label{ii-a-system-and-trust-model}

\textbf{Principals.} Let \(\mathcal{T} = \{t_1, \dots, t_N\}\) be the
set of tenants provisioned on the platform. Each tenant \(t_i\) is
associated with (i) a set of \emph{webhook producers} --- external
services (payment networks, ticketing backends) authorized by \(t_i\) to
emit events --- and (ii) a per-tenant secret \(k_i\) used to
authenticate those events (§IV). The platform operates an \textbf{edge
firewall} \(F\) (a V8 isolate, §III) and an \textbf{origin} \(O\) that
hosts the shared LLM inference service \(\mathcal{M}\). A webhook is an
HTTP request \(r\) carrying a payload \(p\) and metadata, delivered to
\(F\), which either forwards a sanitized request to \(O\) or rejects it.

\textbf{Trust boundaries.} Three boundaries matter. \textbf{B1}, the
public network between producers and \(F\) (untrusted). \textbf{B2},
between \(F\) and \(O\) (mutually authenticated, private). \textbf{B3},
\emph{inside} \(\mathcal{M}\), between the model\textquotesingle s
system/instruction context and the tenant-supplied content --- the
boundary that prompt injection attacks. Classical architectures collapse
B1\textquotesingle s enforcement into \(O\); we relocate it to \(F\),
which shares no persistent state with tenant data.

\textbf{Assets.} (A1) Per-tenant secrets \(k_i\). (A2) Tenant data
resident in or reachable from \(\mathcal{M}\)\textquotesingle s context.
(A3) The integrity of the instruction channel at B3. (A4) Availability
and cost bounds of \(\mathcal{M}\). (A5) The non-repudiable record of
security events.

\textbf{Adversary model.} We assume a \textbf{Dolev--Yao network
attacker} on B1 who can observe, drop, replay, reorder, and inject
arbitrary messages, but cannot break cryptographic primitives. We
additionally assume a \textbf{malicious-but-authenticated tenant}
\(t_a\): a legitimate principal holding a valid \(k_a\) who attempts to
influence another tenant \(t_b\)\textquotesingle s model interactions
(the cross-tenant adversary). We do \emph{not} assume a compromised
\(O\) or a broken HMAC/SHA-256; those are the trust roots. This dual
adversary --- external forger plus insider tenant --- is what
distinguishes multi-tenant LLM webhook security from single-tenant
chatbot hardening.

\subsection{II-B. STRIDE Decomposition with Cross-Boundary
Coupling}\label{ii-b-stride-decomposition-with-cross-boundary-coupling}

STRIDE classifies threats against \emph{data flows}, and its six
categories were defined for systems in which the control plane (code)
and data plane (data) are physically distinct channels. An LLM violates
this premise: instructions and data share one natural-language channel,
so \textbf{Tampering with the data an LLM reads becomes Elevation of
Privilege the instant the model acts on it.} This is not a new taxonomy;
it is an \emph{extended application} of STRIDE in which we retain the
six canonical categories as the primary classification and additionally
annotate the \textbf{coupling chains} (\(T\!\rightarrow\!E\),
\(T\!\rightarrow\!I\)) --- a single event traversing two categories
across a trust boundary --- that a flat per-category classification
would obscure. Table I gives the decomposition over the webhook$\to$LLM
pipeline.

\textbf{TABLE I. STRIDE decomposition of the multi-tenant webhook$\to$LLM
pipeline.}

{\def\LTcaptype{none} % do not increment counter
\begin{table*}[!t]\centering\scriptsize
\resizebox{\textwidth}{!}{%
\begin{tabular}{@{}lllllll@{}}
\toprule
\# & STRIDE & Threat (webhook$\to$LLM) & Asset & Trust bdy & Coupling &
Mitigating layer \\
\midrule



1 & \textbf{S}poofing & Forged/unsigned webhook impersonating a producer
or tenant & A1 & B1 & --- & HMAC tenant binding (§IV) \\
2 & \textbf{T}ampering & Payload mutation in transit & A3 & B1 & --- &
Signature over canonical payload (§IV) \\
3 & \textbf{T}ampering$\to$\textbf{E} & Indirect injection: adversarial
instructions embedded in retrieved/relayed content & A3 & B1,B3 &
\(T\!\rightarrow\!E\) & Heuristic screening (§V) + isolation \\
4 & \textbf{T}ampering$\to$\textbf{I} & Cross-tenant context poisoning:
\(t_a\) mutates shared state later read into \(t_b\)\textquotesingle s
context & A2,A3 & B3 & \(T\!\rightarrow\!I\) & Tenant binding (§IV) +
per-tenant context \\
5 & \textbf{R}epudiation & Attack occurs with no durable attribution &
A5 & B1 & --- & Async threat log (§VI) \\
6 & \textbf{I}nfo. disclosure & System-prompt/context extraction;
cross-tenant leakage & A2 & B3 & --- & Screening (§V) + isolation \\
7 & \textbf{D}oS & Token flood / CPU / subrequest exhaustion;
"denial-of-wallet" cost amplification & A4 & B1 & --- & Edge rejection
within CPU budget (§III,§V) \\
8 & \textbf{E}levation & Direct injection/jailbreak executes with the
agent\textquotesingle s tool/data privileges & A2,A3 & B3 & --- &
Screening (§V) + least-privilege tools \\
\bottomrule
\end{tabular}}
\end{table*}
}

\textbf{The lateral (multi-tenancy) axis.} Orthogonal to STRIDE, each
threat has a \emph{same-tenant} and a \emph{cross-tenant} manifestation.
Same-tenant injection (a tenant attacking its own model surface) is
comparatively benign and well-studied; the platform\textquotesingle s
obligation is the \textbf{cross-tenant} cases --- rows 4 and 6 --- where
one tenant\textquotesingle s input compromises another\textquotesingle s
confidentiality or integrity. We therefore treat the effective threat
space as \(\text{STRIDE} \times \{\text{same},
\text{cross}\}\) and prioritize the cross-tenant quadrant. This is the
precise sense in which our contribution is \emph{multi-tenant} security
rather than generic LLM hardening.

\textbf{Why the coupling annotation matters for defense placement.}
Because \(T\!\rightarrow\!E\) (row 3) originates as \emph{Tampering on
B1} but detonates as \emph{Elevation at B3}, it can be intercepted at
\emph{either} boundary. Origin-only defenses can act only at B3, after
the model has already ingested the payload. Relocating screening to
\(F\) lets us break the chain at B1 --- before the data ever reaches the
channel where it can become an instruction. This is the architectural
justification, in threat-model terms, for the edge relocation argued
informally in §I-B.

\subsection{II-C. Prompt-Injection Taxonomy (OWASP
LLM01:2025)}\label{ii-c-prompt-injection-taxonomy-owasp-llm012025}

We adopt the taxonomy of OWASP\textquotesingle s \emph{Top 10 for LLM
Applications}, entry \textbf{LLM01: Prompt Injection}, as the
authoritative reference {[}OWASP-LLM01{]}, corroborated in the
literature {[}arXiv-2410.21146{]}. Prompt injection partitions into:

\begin{itemize}
\tightlist
\item
  \textbf{Direct injection.} The attacker supplies malicious
  instructions through inputs the model processes in real time, aiming
  to override the intended system instructions (e.g., "ignore previous
  instructions and \ldots"). In the webhook setting, the payload body is
  the injection vector.
\item
  \textbf{Indirect injection.} Adversarial instructions are embedded in
  \emph{externally retrieved or relayed} content --- a document, a web
  page, an upstream event field --- that the model consumes as data but
  interprets as instruction. This is the vector for row 3 of Table I and
  the more insidious because the injecting party need not be the
  requesting party.

  \begin{itemize}
  \tightlist
  \item
    \textbf{Stored injection} is a subcategory of indirect in which the
    payload is persisted (e.g., in a record or knowledge base) and
    executes on a \emph{later}, possibly different-tenant, retrieval ---
    the mechanism underlying cross-tenant context poisoning (row 4).
  \end{itemize}
\end{itemize}

We further sub-classify injection \emph{techniques} (jailbreak families)
as an implementation concern of the screening layer ---
instruction-override, role-play/persona ("DAN"-style), system-prompt
extraction, delimiter/tag spoofing, and control-token injection --- and
enumerate their signatures in §V and Appendix C. This taxonomy is
\emph{verified}: the direct/indirect (with stored as an indirect
subcategory) partition and its OWASP grounding were confirmed at high
confidence.

\subsection{II-D. Formalizing Cross-Tenant Data
Leakage}\label{ii-d-formalizing-cross-tenant-data-leakage}

Let \(c_b\) denote the model context assembled for a request on behalf
of tenant \(t_b\): it comprises a platform system prompt \(s\),
tenant-scoped data \(d_b\), and the request payload \(p\).
\textbf{Cross-tenant leakage} is any execution in which information
derived from \(d_{a}\) (\(a \neq b\)) becomes observable in the response
to \(t_b\), or vice versa. Two failure modes generate it.

\begin{enumerate}
\def\labelenumi{\arabic{enumi}.}
\tightlist
\item
  \textbf{Shared-context leakage.} If context assembly places \(d_a\)
  and \(d_b\) in the same \(c\) (e.g., a shared conversation buffer, a
  mis-scoped retrieval, a global cache), an extraction injection (row 6)
  in \(p\) can exfiltrate \(d_a\). Formally, isolation requires that for
  every request on behalf of \(t_b\), \(c_b \cap d_a = \emptyset\) for
  all \(a \neq b\) --- a property the \emph{architecture}, not the
  model, must guarantee.
\item
  \textbf{Shared-key attribution failure.} If a single secret \(k\)
  authenticates all tenants, a valid signature proves only that
  \emph{some} authorized party sent the event, not \emph{which} tenant.
  An attacker (or a confused-deputy producer) can then present content
  that the origin attributes to the wrong tenant, causing \(p\) authored
  under \(t_a\)\textquotesingle s authority to execute in \(c_b\).
  Preventing this is exactly the tenant-\emph{binding} property we
  construct and prove in §IV: the signature must authenticate not just
  \emph{authenticity} but \emph{tenant identity}.
\end{enumerate}

These two modes motivate the two independent controls the architecture
provides --- per-tenant context isolation at the origin (assumed, and
out of scope for the edge firewall) and per-tenant cryptographic binding
at the edge (§IV, our contribution). The threat model thus reduces the
cross-tenant obligation to a property we can state and prove
cryptographically, which §IV does.

\begin{center}\rule{0.5\linewidth}{0.5pt}\end{center}

\subsubsection{Citation keys}\label{citation-keys}

\begin{itemize}
\tightlist
\item
  \textbf{{[}OWASP-LLM01{]}} OWASP, "LLM01:2025 Prompt Injection," Top
  10 for LLM Applications,
  genai.owasp.org/llmrisk/llm01-prompt-injection.
\item
  \textbf{{[}arXiv-2410.21146{]}} "\emph{(indirect/stored injection
  taxonomy)}," arXiv:2410.21146.
\end{itemize}

\begin{quote}
decomposition and cross-tenant formalization in §II-A/B/D are original
analytical framing built on the verified taxonomy; they assert no
unverified external facts.
\end{quote}

\section{III. Edge Interception Layer: The V8 Isolate Execution
Envelope}\label{iii-edge-interception-layer-the-v8-isolate-execution-envelope}

This section characterizes the runtime in which the firewall executes
and converts its published resource limits into an \emph{operation
budget} against which §V\textquotesingle s algorithms are checked. We
first describe the isolate execution model (§III-A), fix the resource
envelope from primary sources (§III-B), derive the cost model and budget
(§III-C), restate the CPU/wall-clock correction and its consequences
(§III-D), and give the reverse-proxy dataflow (§III-E).

\subsection{III-A. The Isolate Execution
Model}\label{iii-a-the-isolate-execution-model}

Cloudflare Workers execute in \textbf{V8 isolates} rather than
per-request containers or processes. An isolate is a lightweight,
memory-safe sandbox within a shared V8 runtime; thousands coexist in one
OS process, and a single isolate serves \emph{many} concurrent requests.
This has three consequences that shape the firewall\textquotesingle s
design.

\begin{enumerate}
\def\labelenumi{\arabic{enumi}.}
\tightlist
\item
  \textbf{No per-request process spin-up.} Unlike a
  container-per-request model, isolate reuse means module-scope
  initialization (parsing configuration, \emph{building the
  Aho--Corasick automaton}) executes once when the isolate is created
  and is then amortized across every request that isolate serves. Only
  work performed \emph{inside} the request handler counts against
  per-request CPU.
\item
  \textbf{Cold start vs. warm path.} A request routed to a location with
  no warm isolate pays a one-time initialization cost (the module
  top-level), after which the isolate stays warm and subsequent requests
  skip it. Our design places all heavy precomputation at module top
  level precisely so that it is paid at most once per isolate lifetime,
  never on the warm request path.
\item
  \textbf{Single-threaded, event-loop concurrency.} Each isolate runs
  JavaScript on a single thread; concurrency is cooperative via the
  event loop. CPU-bound work therefore \emph{blocks} the isolate for its
  duration, which is why the per-request CPU cost of screening must be
  small in absolute terms, not merely asymptotically linear. This
  motivates the operation-budget analysis of §III-C.
\end{enumerate}

The isolate model also underlies the memory constraint: because one
isolate hosts many requests, the 128 MB cap (§III-B) is a
\emph{per-isolate}, not per-request, limit, and long-lived module-scope
structures (the automaton, signature tables) are counted against it
once, not per request.

\subsection{III-B. The Resource
Envelope}\label{iii-b-the-resource-envelope}

Table II fixes the operative limits from primary Cloudflare
documentation. \textbf{Every figure is date-stamped (accessed
2026-07-04); the platform revises these quarterly, and any reuse of this
table must re-verify.}

\textbf{TABLE II. Cloudflare Workers resource envelope (accessed
2026-07-04).}

{\def\LTcaptype{none} % do not increment counter
\begin{table*}[!t]\centering\scriptsize
\resizebox{\textwidth}{!}{%
\begin{tabular}{@{}llll@{}}
\toprule
Resource & Free plan & Paid (Standard) plan & Source \\
\midrule



CPU time / invocation & \textbf{10 ms} & \textbf{30 s default},
configurable to \textbf{300 s (5 min)} via
\passthrough{\lstinline!limits.cpu\_ms!} & {[}C-limits{]},
{[}C-changelog-2025{]} \\
Wall-clock duration (HTTP trigger) & no hard limit & no hard limit / no
charge & {[}C-limits{]}, {[}C-pricing{]} \\
Memory / isolate & 128 MB & 128 MB & {[}C-limits{]} \\
Subrequests / invocation & 50 external + 1,000 to CF services &
\textbf{10,000 default} (to 10 M), since 2026-02-11 & {[}C-limits{]},
{[}C-changelog-2026{]} \\
I/O wait counted as CPU? & \textbf{No} & \textbf{No} & {[}C-limits{]} \\
\bottomrule
\end{tabular}}
\end{table*}
}

Three properties of this envelope are decisive for the architecture:

\begin{itemize}
\tightlist
\item
  \textbf{CPU time excludes I/O wait.} Time awaiting the
  \passthrough{\lstinline!fetch!} to the origin LLM, a KV read for a
  nonce (§IV), or the database write (§VI) does \emph{not} accrue
  against the CPU limit {[}C-limits{]}. Consequently the
  firewall\textquotesingle s CPU budget is consumed \emph{only} by local
  computation --- signature scanning, HMAC verification, sanitization
  --- and not by the network operations that dominate wall-clock time.
\item
  \textbf{Wall-clock is effectively unbounded for HTTP-triggered
  Workers.} The webhook path is HTTP-triggered, so there is no duration
  cap to design against; the "\textless50 ms" target (§VII) is a
  self-imposed \emph{latency SLO}, not a platform limit.
\item
  \textbf{The subrequest budget is ample.} Forwarding to origin plus at
  most a small constant number of KV/queue operations per request sits
  far within even the Free-plan 50-external ceiling.
\end{itemize}

\subsection{III-C. From CPU Limit to Operation
Budget}\label{iii-c-from-cpu-limit-to-operation-budget}

Asymptotic complexity bounds \emph{scaling} but not \emph{absolute}
latency; a linear algorithm with a large constant on an unbounded input
can exceed any fixed budget. A defensible sub-millisecond guarantee
therefore requires three ingredients: (i) a hard bound on input size,
(ii) a per-operation constant on the target runtime, and (iii) a
worst-case operation count. We supply (i) and (iii) here analytically
and pin (ii) empirically in §VII.

\textbf{Operation budget.} Let \(\rho\) (operations $\cdot$ ms\(^{-1}\)) be
the effective scalar-operation throughput of the edge V8 runtime for the
byte-scanning workload of §V, and let \(L_{\text{CPU}}\) be the
per-invocation CPU limit. The available operation budget for one request
is \fiteq{B = \rho \cdot L_{\text{CPU}}.} On the Free plan,
\(L_{\text{CPU}} = 10\text{ ms}\); on paid plans the \emph{default} is
\(L_{\text{CPU}} =
30{,}000\text{ ms}\). We deliberately evaluate the firewall against the
\textbf{Free-plan} budget as the worst case --- if screening fits in 10
ms of CPU, it fits everywhere. \(\rho\) is left symbolic here and
measured in §VII; the analysis below yields a bound of the form
"\(C_{\text{req}}/B\)" that is independent of \(\rho\)\textquotesingle s
exact value once \(\rho\) is known.

\textbf{Input bound.} The firewall rejects any request whose body
exceeds a configured maximum \(N_{\max}\) bytes \emph{before} invoking
the scanner (Algorithm in §V). Oversized-payload rejection is
simultaneously (a) a precondition for the latency guarantee and (b) a
DoS control (Table I, row 7). We take \(N_{\max}\) as a deployment
parameter (default 128 KiB in our implementation, §V/Phase 2).

\textbf{Per-request cost.} Under the isolate model (§III-A), the
O(\(m\)) Aho--Corasick construction over total signature length \(m\)
executes once at module init and contributes \textbf{zero} to the
per-request budget. The per-request work is therefore:
\fiteq{C_{\text{req}} \;\le\; \underbrace{c_{\text{ac}}\,(N_{\max} + z)}_{\text{signature scan (§V-B)}}
\;+\; \underbrace{k\,c_{\text{lin}}\,N_{\max}}_{k \text{ linear feature passes (§V-C)}}
\;+\; \underbrace{c_{\text{hmac}}\,N_{\max}}_{\text{HMAC over payload (§IV)}},}
where \(z \le N_{\max}\) is the number of signature matches (bounded by
input length and further capped by early-exit, §V-B), \(k\) is the fixed
number of linear feature passes (entropy, structural --- a small
constant, \(k \le 4\)), and \(c_{\bullet}\) are per-byte constants.
Since \(z \le N_{\max}\), this simplifies to
\fiteq{C_{\text{req}} \;\le\; \big(2c_{\text{ac}} + k\,c_{\text{lin}} + c_{\text{hmac}}\big)\,N_{\max}
\;=\; \kappa\, N_{\max},} a \textbf{constant \(\kappa\) times a bounded
input} --- the only form that supports an absolute latency claim. The
sub-millisecond guarantee is then the assertion
\(\kappa N_{\max} \ll B\), i.e.
\fiteq{\frac{C_{\text{req}}}{B} = \frac{\kappa N_{\max}}{\rho\,L_{\text{CPU}}} \ll 1,}
which §VII establishes by measuring \(\kappa/\rho\) (the wall-normalized
per-byte cost) and evaluating at \(N_{\max} = 128\text{ KiB}\),
\(L_{\text{CPU}} = 10\text{ ms}\). The critical structural point,
provable \emph{without} the constant, is that per-request cost is
\textbf{strictly linear in a bounded input with all superlinear and
one-time work excluded} --- there is no per-request construction, no
backtracking (Aho--Corasick is backtrack-free), and no unbounded loop.

\subsection{III-D. The CPU/Wall-Clock Correction and Its
Consequences}\label{iii-d-the-cpuwall-clock-correction-and-its-consequences}

As established in §I-C, the commonly cited "50 ms Worker CPU limit" is a
\textbf{legacy artifact} of the deprecated Bundled usage model,
auto-applied during the 2024-03-01 migration to Standard pricing
{[}C-pricing{]}; it is not the contemporary envelope. Table II
supersedes it. Two engineering consequences follow directly and are
exploited by the architecture:

\begin{enumerate}
\def\labelenumi{\arabic{enumi}.}
\tightlist
\item
  \textbf{The screening layer runs with vast headroom, not at the
  margin.} Against a 10 ms Free-plan CPU ceiling --- and 30,000 ms on
  paid plans --- a \(\kappa N_{\max}\) cost on a 128 KiB bound is orders
  of magnitude under budget (§VII quantifies the ratio). The design
  problem is not "fit inside a tight CPU limit" but "spend a tiny,
  \emph{bounded} fraction of an ample budget so that screening is
  invisible in the wall-clock latency envelope."
\item
  \textbf{I/O is free against the CPU budget, so decoupling is natural.}
  Because awaiting the origin \passthrough{\lstinline!fetch!} and the
  threat-log write does not accrue CPU (§III-B), the firewall can
  forward the request and hand off logging (§VI) without those
  operations competing with screening for the CPU limit. The CPU limit
  constrains \emph{only} the local inspection, which §III-C has bounded.
\end{enumerate}

\subsection{III-E. Reverse-Proxy
Dataflow}\label{iii-e-reverse-proxy-dataflow}

The firewall is a reverse proxy on the request path (Fig. 1). For each
inbound webhook:

\begin{lstlisting}
                         +------------------------ Edge Isolate F -----------------------+
 Producer --HTTP(B1)-->  |  (1) size guard: reject if |body| > N_max        [O(1)]        |
                         |  (2) HMAC-SHA256 verify + tenant binding (§IV)   [c_hmac$\cdot$N]     |
                         |  (3) canonicalize + sanitize input (§V-E)        [O(N)]         |
                         |  (4) Aho–Corasick signature scan (§V-B)          [c_ac$\cdot$(N+z)]   |
                         |  (5) structural + entropy scoring (§V-C)         [k$\cdot$c_lin$\cdot$N]    |
                         |  (6) decision: forward | reject                  [O(1)]         |
                         |        | forward (B2)                                            |
                         |        v                                                         |
                         |   fetch $\to$ Origin O / model M   -- I/O, not CPU --> response      |
                         |        |                                                         |
                         |  (7) ctx.waitUntil(logThreat(...))  -- async, post-response -->  |--> PostgreSQL (§VI)
                         +---------------------------------------------------------------+
\end{lstlisting}

\textbf{Fig. 1.} Request lifecycle in the edge isolate. Steps (1)--(6)
are the CPU-bounded critical path analyzed in §III-C; the origin
\passthrough{\lstinline!fetch!} and the
\passthrough{\lstinline!waitUntil!}-deferred threat log (§VI) are I/O
and do not accrue against the CPU limit. Steps (2), (4), (5) are the
substance of §IV and §V; the size guard (1) supplies the \(N_{\max}\)
bound that makes the latency analysis absolute.

The ordering is security-critical: authentication (step 2) precedes
screening (steps 3--5), so unauthenticated traffic is rejected before
any content-inspection cost is incurred, and the size guard (step 1)
precedes everything, so no unbounded input reaches any linear pass.

\begin{center}\rule{0.5\linewidth}{0.5pt}\end{center}

\subsubsection{Citation keys}\label{citation-keys-1}

\begin{itemize}
\tightlist
\item
  \textbf{{[}C-limits{]}} Cloudflare, "Workers --- Limits,"
  developers.cloudflare.com/workers/platform/limits (accessed
  2026-07-04).
\item
  \textbf{{[}C-pricing{]}} Cloudflare, "Workers --- Pricing,"
  developers.cloudflare.com/workers/platform/pricing (accessed
  2026-07-04).
\item
  \textbf{{[}C-changelog-2025{]}} Cloudflare, "Higher CPU limits for
  Workers," changelog, 2025-03-25.
\item
  \textbf{{[}C-changelog-2026{]}} Cloudflare, "Increased subrequest
  limits," changelog, 2026-02-11.
\end{itemize}

\begin{quote}
3-0). The cost model (§III-C) is original analysis; the constant
\(\rho\) (and hence \(\kappa/\rho\)) is deferred to empirical
measurement in §VII and asserted nowhere as a literature value.
\end{quote}

\section{IV. Cryptographic Tenant
Isolation}\label{iv-cryptographic-tenant-isolation}

Section II-D reduced the cross-tenant obligation to a cryptographic
property: a webhook must authenticate not merely \emph{authenticity}
("some authorized party sent this") but \emph{tenant identity} ("tenant
\(t_b\) authorized this"). This section discharges that obligation. We
fix MAC preliminaries (§IV-A), define the tenant-bound signing scheme
(§IV-B), prove tenant binding for static keys (§IV-C, Theorem 1), lift
the result to key rotation (§IV-D, Theorem 2), and treat replay as an
orthogonal freshness property (§IV-E). §IV-F states why the tenant
identifier must reside \emph{inside} the signed message.

\subsection{IV-A. Preliminaries: MACs, EUF-CMA, and
HMAC}\label{iv-a-preliminaries-macs-euf-cma-and-hmac}

\textbf{MAC syntax.} A message authentication code is a triple
\(\Pi=(\mathsf{KGen},\mathsf{Mac},
\mathsf{Vrfy})\): \(\mathsf{KGen}\) outputs a key \(k\in\{0,1\}^n\);
\(\mathsf{Mac}_k(m)\to\tau\) produces a tag;
\(\mathsf{Vrfy}_k(m,\tau)\in\{0,1\}\) verifies. Correctness requires
\(\mathsf{Vrfy}_k(m,\mathsf{Mac}_k(m))=1\) for all \(k,m\).

\textbf{EUF-CMA.} Existential unforgeability under chosen-message attack
is defined by the game
\(\mathbf{Exp}^{\text{euf-cma}}_{\Pi}(\mathcal{A})\):

\begin{enumerate}
\def\labelenumi{\arabic{enumi}.}
\tightlist
\item
  Challenger runs \(k\leftarrow\mathsf{KGen}(1^n)\) and initializes
  \(\mathcal{Q}\leftarrow\emptyset\).
\item
  \(\mathcal{A}\) is given oracle access to \(\mathsf{Mac}_k(\cdot)\)
  (each query \(m\) appended to \(\mathcal{Q}\)) and
  \(\mathsf{Vrfy}_k(\cdot,\cdot)\).
\item
  \(\mathcal{A}\) outputs \((m^*,\tau^*)\).
\item
  \(\mathcal{A}\) \textbf{wins} iff
  \(\mathsf{Vrfy}_k(m^*,\tau^*)=1 \wedge m^*\notin\mathcal{Q}\).
\end{enumerate}

The advantage is
\(\mathbf{Adv}^{\text{euf-cma}}_{\Pi}(\mathcal{A})=\Pr[\mathcal{A}\text{ wins}]\),
and \(\Pi\) is EUF-CMA-secure if this is negligible for all PPT
\(\mathcal{A}\). \textbf{The freshness condition
\(m^*\notin\mathcal{Q}\) is essential and is precisely why
unforgeability says nothing about replay of an already-signed message}
(§IV-E).

\textbf{HMAC.} For a hash \(H\) with block length \(B\),
\(\mathsf{HMAC}_k(m)=H\big((k'\oplus\mathrm{opad})\,\|\,
H((k'\oplus\mathrm{ipad})\,\|\,m)\big)\), with
\(\mathrm{ipad}=\texttt{0x36}^B\), \(\mathrm{opad}=\texttt{0x5C}^B\),
and \(k'\) the key zero-padded to \(B\) bytes {[}RFC2104{]},
{[}FIPS198-1{]}. We rely on the standard result that HMAC is a
pseudorandom function (PRF) when the underlying compression function is
a PRF {[}BCK96{]}, and that any PRF is an EUF-CMA-secure MAC with
\[\mathbf{Adv}^{\text{euf-cma}}_{\mathsf{HMAC}}(\mathcal{A}) \;\le\;
\mathbf{Adv}^{\text{prf}}_{\mathsf{HMAC}}(\mathcal{B}) + 2^{-n},\tag{1}\]
where \(n\) is the tag length in bits (for HMAC-SHA256, \(n=256\), so
\(2^{-n}\) is cryptographically negligible). We treat (1) as a trust
root (§II-A) and reduce all subsequent claims to it.

\textbf{Multi-user security.} Because a multi-tenant deployment
instantiates \(u\) \emph{independent} keys (one per tenant, or per
tenant-epoch under rotation), the operative notion is \emph{multi-user}
security. In the multi-user PRF (mu-PRF) game, \(\mathcal{A}\) interacts
with \(u\) independent instances and distinguishes them jointly from
\(u\) random functions; in the multi-user EUF-CMA (mu-EUF-CMA) game,
\(\mathcal{A}\) wins by producing a forgery under \emph{any} one of the
\(u\) instances. The naive hybrid bound loses a factor \(u\),
\(\mathbf{Adv}^{\text{mu-prf}}_{\mathsf{HMAC},u}\le u\cdot\mathbf{Adv}^{\text{prf}}_{\mathsf{HMAC}}\),
which is unacceptable when \(u\) scales to \(10^5\)--\(10^6\) tenants.
We instead invoke the \emph{tight} multi-user analysis of HMAC
{[}BBT16{]}, under which
\[\mathbf{Adv}^{\text{mu-euf-cma}}_{\mathsf{HMAC},u}(\mathcal{A})\;\le\;
\mathbf{Adv}^{\text{mu-prf}}_{\mathsf{HMAC},u}(\mathcal{A})+Q_v\,2^{-n},\tag{2}\]
where the mu-PRF term \textbf{carries no linear-in-\(u\) factor} --- it
is bounded by the adversary\textquotesingle s \emph{aggregate} query
budget and a birthday term, independent of the number of instances ---
and \(Q_v\) is the number of verification queries. Equations (1)--(2)
are the only cryptographic assumptions used below.

\subsection{IV-B. The Tenant-Bound Webhook Signing
Scheme}\label{iv-b-the-tenant-bound-webhook-signing-scheme}

For tenant \(t_i\) holding key \(k_{i,\kappa}\) under key-id \(\kappa\),
define the \textbf{canonical signed message}
\fiteq{m \;=\; \mathsf{tid}_i \,\|\, \kappa \,\|\, t_s \,\|\, \eta \,\|\, H_{\text{body}},}
where \(\mathsf{tid}_i\) is the tenant identifier, \(\kappa\) the
key-id, \(t_s\) a Unix timestamp, \(\eta\) a random nonce (\(\lambda\)
bits), and \(H_{\text{body}}=\mathsf{SHA256}(\text{payload})\) a digest
binding the body. The transmitted signature is
\(\tau=\mathsf{HMAC}_{k_{i,\kappa}}(m)\), sent with
\((\mathsf{tid}_i,\kappa,t_s,\eta)\) in headers. Each field is
length-prefixed (or delimited by a byte absent from the field alphabet)
so that the encoding is injective --- no two distinct field tuples share
a serialization, foreclosing canonicalization-ambiguity attacks. This
mirrors the Stripe scheme (signed payload \(=\)
\passthrough{\lstinline!timestamp . "." . body!},
\passthrough{\lstinline!v1=HMAC-SHA256!}) {[}Stripe-sig{]} but
additionally binds \(\mathsf{tid}_i\) and \(\kappa\).

Verification at the edge (Fig. 1, step 2): recompute
\(\tau'=\mathsf{HMAC}_{k_{i,\kappa}}(m)\) using the key selected by
\((\mathsf{tid}_i,\kappa)\) and accept iff \(\tau'=\tau\) under a
constant-time comparison (to avoid timing side channels), the timestamp
is fresh, and the nonce is unseen (§IV-E).

\subsection{IV-C. Tenant Binding under Static Keys (Theorem
1)}\label{iv-c-tenant-binding-under-static-keys-theorem-1}

\textbf{Tenant-binding game} \(\mathbf{Exp}^{\text{bind}}\). Let each
tenant \(t_i\in\mathcal{T}\) have an independent key
\(k_i\leftarrow\mathsf{KGen}(1^n)\) (single key per tenant in this
subsection). The adversary \(\mathcal{A}\):

\begin{enumerate}
\def\labelenumi{\arabic{enumi}.}
\tightlist
\item
  selects a corrupt set \(\mathcal{C}\subset\mathcal{T}\) and receives
  \(\{k_i : t_i\in\mathcal{C}\}\);
\item
  for every uncorrupted tenant \(t_j\notin\mathcal{C}\), obtains oracle
  access to \(\mathsf{Mac}_{k_j}(\cdot)\) (this over-approximates the
  Dolev--Yao capability: observing \(t_j\)\textquotesingle s legitimate
  webhooks is a \emph{known}-message attack, a special case of
  chosen-message);
\item
  outputs a target \(t_b\notin\mathcal{C}\) and a pair
  \((m^*,\tau^*)\) with \(\mathsf{tid}(m^*)=
  \mathsf{tid}_b\).
\end{enumerate}

\(\mathcal{A}\) \textbf{wins} iff
\(\mathsf{Vrfy}_{k_b}(m^*,\tau^*)=1\) and \(m^*\) was never returned
by \(t_b\)\textquotesingle s signing oracle. Intuitively: the attacker,
even owning every other tenant\textquotesingle s key and seeing all of
\(t_b\)\textquotesingle s traffic, produces a \emph{new} webhook that
the system attributes to \(t_b\).

\begin{quote}
\textbf{Theorem 1 (Tenant binding, static keys --- multi-user).} Let
\(u=|\mathcal{T}\setminus
\mathcal{C}|\) be the number of uncorrupted tenants. For the scheme of
§IV-B with independent per-tenant keys and any PPT adversary
\(\mathcal{A}\), there exists a PPT \(\mathcal{B}\) with
\fiteq{\mathbf{Adv}^{\text{bind}}(\mathcal{A}) \;\le\;
\mathbf{Adv}^{\text{mu-euf-cma}}_{\mathsf{HMAC},u}(\mathcal{B})
\;\le\; \mathbf{Adv}^{\text{mu-prf}}_{\mathsf{HMAC},u}(\mathcal{B})+Q_v\,2^{-n}.}
By the tight multi-user security of HMAC (Eq. 2, {[}BBT16{]}) the
right-hand side is \textbf{independent of the tenant-pool size \(u\)}
--- it is bounded by \(\mathcal{A}\)\textquotesingle s aggregate query
budget and a birthday term, not by the number of tenants.
\end{quote}

\emph{Proof (multi-user reduction, no target guessing).} \(\mathcal{B}\)
plays mu-EUF-CMA against the \(u\) uncorrupted-tenant instances
\(\{k_j : t_j\notin\mathcal{C}\}\). It generates the corrupt
tenants\textquotesingle{} keys itself and answers their key-reveal and
signing queries directly; each signing query for an uncorrupted \(t_j\)
is forwarded to instance \(j\)\textquotesingle s \(\mathsf{Mac}\)
oracle, recording the queried \(m\). When \(\mathcal{A}\) halts with
\((m^*,\tau^*)\) where \(\mathsf{tid}(m^*)=\mathsf{tid}_b\) for some
uncorrupted \(t_b\) and \(m^*\) was never signed by \(t_b\), then ---
because the encoding is injective and \(\mathsf{tid}\) occupies a fixed
field --- \((b, m^*,\tau^*)\) is verbatim a valid forgery under
instance \(b\) with \(m^*\) fresh for that instance: a win in the
mu-EUF-CMA game. \textbf{No guess of the target is required}, because a
forgery under \emph{any} uncorrupted instance already wins the
multi-user game; this is exactly what removes the factor-\(u\) loss. The
second inequality is the generic mu-PRF \(\Rightarrow\) mu-MAC step (Eq.
2). \(\square\)

The theorem formalizes the §II-D requirement: because \(\mathsf{tid}_i\)
is \emph{inside} the MAC input, attributing a webhook to \(t_b\) is
exactly as hard as forging HMAC in the multi-user game --- independent
of how many \emph{other} tenant keys the adversary holds \emph{and}
independent of the tenant population. This is the shared-key attribution
failure (§II-D, mode 2) provably eliminated, with a bound that does not
degrade at enterprise scale.

\subsection{IV-D. Key Rotation in the Formal Model (Theorem
2)}\label{iv-d-key-rotation-in-the-formal-model-theorem-2}

Static keys are unrealistic: keys must rotate for compromise recovery
and hygiene. We model rotation without abandoning Theorem 1.

\textbf{Rotation model.} Each tenant maintains a set of keys
\(\{k_{i,\kappa}\}\) indexed by key-id \(\kappa\), each with a validity
window \([\,a_\kappa, b_\kappa\,)\). Windows \emph{overlap}: when
rotating from \(\kappa\) to \(\kappa+1\), both are valid for a rollover
interval, so in-flight producers signing under the old key are not
rejected. Let \(L=\max_i \max_t |\{\kappa : t\in[a_\kappa,b_\kappa)\}|\)
be the maximum number of simultaneously valid keys for any tenant at any
time (typically \(L=2\)). The verifier selects the key by the
\emph{authenticated} \(\kappa\) carried in \(m\) and accepts iff the tag
validates under \(k_{i,\kappa}\) and \(\kappa\) is currently valid.

\begin{quote}
\textbf{Theorem 2 (Tenant binding under rotation).} Let
\(u'=\sum_{t_i\notin\mathcal{C}}
|\{\kappa : k_{i,\kappa}\text{ currently valid}\}| \le uL\) be the total
number of simultaneously-valid uncorrupted keys. Then
\fiteq{\mathbf{Adv}^{\text{bind-rot}}(\mathcal{A}) \;\le\;
\mathbf{Adv}^{\text{mu-euf-cma}}_{\mathsf{HMAC},u'}(\mathcal{B})
\;\le\; \mathbf{Adv}^{\text{mu-prf}}_{\mathsf{HMAC},u'}(\mathcal{B})+Q_v\,2^{-n},}
which by Eq. 2 is \textbf{independent of both \(u\) and the overlap
multiplicity \(L\)}.
\end{quote}

\emph{Proof (each valid key is an instance).} Treat every
currently-valid uncorrupted pair \((t_i,\kappa)\) as one of \(u'\)
independent mu-EUF-CMA instances. The rotation verifier accepts a
\(t_b\)-attributed message iff its tag validates under \emph{some} valid
\((t_b,\kappa)\) --- i.e., iff it is a forgery under one of those
instances. The reduction of Theorem 1 applies verbatim with the instance
set enlarged from \(u\) to \(u'\): a win for \(\mathcal{A}\) is a
forgery under some instance, winning the mu-EUF-CMA game with no target
or key guess. Because the tight multi-user bound (Eq. 2) has no
linear-in-instance-count factor, enlarging \(u\to u'\le uL\) does not
degrade the bound. \(\square\)

\textbf{Consequences and strategy.} (i) Under the tight multi-user bound
the security is independent of the overlap multiplicity \(L\); \(L\)
affects only \emph{operational} exposure, not the reduction. Keeping the
rollover interval short still matters because it bounds the
\emph{compromise window} (iii), not the advantage. (Under the weaker
generic bound one would pay a factor \(u'\le uL\); we avoid this via
{[}BBT16{]}.) (ii) Because \(\kappa\) is authenticated inside \(m\), an
attacker cannot force verification under a \emph{revoked} key (a
downgrade): tampering with \(\kappa\) breaks the tag. (iii) On
compromise of \(k_{i,\kappa}\), the operator advances \(\kappa\) and
shrinks the old window to \(\{\)now\(\}\); the exposure is bounded by
the rollover interval. Keys are stored in the edge secret store / KV and
selected by \((\mathsf{tid}_i,\kappa)\) at verification; retrieval is
I/O, not CPU (§III-B). (iv) Forward secrecy is \emph{not} claimed: HMAC
keys are symmetric, so a leaked \(k_{i,\kappa}\) forges messages within
\(\kappa\)\textquotesingle s window. Rotation bounds, but does not
retroactively protect, that window.

\subsection{IV-E. Replay Resistance as an Orthogonal Freshness
Property}\label{iv-e-replay-resistance-as-an-orthogonal-freshness-property}

Theorem 1\textquotesingle s freshness clause (\(m^*\notin\mathcal{Q}\))
means a \emph{replayed} legitimate webhook \((m,\tau)\) --- where \(m\)
\emph{was} signed by \(t_b\) --- is \textbf{not} an unforgeability
break. Replay is a distinct property requiring freshness, which the
\(t_s\) and \(\eta\) fields provide.

\textbf{Freshness game.} The verifier maintains a nonce cache
\(\mathcal{N}\) and accepts \((m,\tau)\) only if: (a) \(\tau\) verifies
(§IV-C); (b) \(|t_{\text{now}}-t_s|\le\Delta\) (timestamp tolerance);
and (c) \(\eta\notin\mathcal{N}\), after which \(\eta\) is inserted with
TTL \(2\Delta\). An adversary replaying a captured \((m,\tau)\) succeeds
only if it arrives (i) within \(\Delta\) of \(t_s\) \emph{and} (ii) with
\(\eta\) already evicted from \(\mathcal{N}\). But \(\mathcal{N}\)
retains every nonce for \(2\Delta\ge\Delta\), so within the timestamp
window the nonce is necessarily still present and the replay is
rejected; outside the window the timestamp check rejects it. Hence
\fiteq{\Pr[\text{replay accepted}] \;\le\; \Pr[\text{nonce-store loss within }2\Delta]\;+\;2^{-\lambda},}
where the first term is the probability the nonce store fails to retain
\(\eta\) over the window (an availability parameter of the KV/DO store)
and \(2^{-\lambda}\) bounds an accidental nonce collision causing false
eviction. With \(\lambda=128\) the collision term is negligible, so
replay resistance reduces to nonce-store reliability over \(2\Delta\).

\textbf{Parameterization and cost.} \(\Delta\) trades replay window
against tolerance to clock skew and producer/delivery latency; Stripe
uses \(\Delta=300\text{ s}\) by default {[}Stripe-sig{]}, which we adopt
as a baseline. The nonce check is a single keyed lookup in edge KV or a
Durable Object; it is a subrequest (I/O), so it does \textbf{not}
consume the CPU budget of §III-C --- replay defense is free against the
sub-millisecond screening analysis. A stateless fallback
(timestamp-only, no nonce) degrades to "at-most-one replay per
\(\Delta\) window per message" and is offered as a lower-assurance mode
when KV latency is unacceptable.

\subsection{IV-F. Why the Tenant Identifier Must Be Inside the
Signature}\label{iv-f-why-the-tenant-identifier-must-be-inside-the-signature}

If \(\mathsf{tid}_i\) were transmitted only as an \emph{unauthenticated}
header (outside \(m\)), an attacker holding any single valid
\((m,\tau)\) under key \(k\) could attach an arbitrary \(\mathsf{tid}\)
header, and --- if the origin selected the verification key by the
\emph{header} rather than by the signed field --- cause \(t_b\)-authored
content to execute in \(t_a\)\textquotesingle s context or vice versa.
This is exactly the confused-deputy instantiation of §II-D mode 2.
Binding \(\mathsf{tid}_i\) (and \(\kappa\)) \emph{inside} the MAC input
makes the identifier immutable under the unforgeability guarantee: any
change to \(\mathsf{tid}\) changes \(m\), which invalidates \(\tau\)
unless the attacker can forge --- contradicting Theorem 1. The binding
is therefore not a convention but a proven property, which is the sense
in which §II-D\textquotesingle s isolation obligation is
\emph{discharged} rather than merely \emph{addressed}.

\begin{center}\rule{0.5\linewidth}{0.5pt}\end{center}

\subsubsection{Citation keys}\label{citation-keys-2}

\begin{itemize}
\tightlist
\item
  \textbf{{[}RFC2104{]}} H. Krawczyk, M. Bellare, R. Canetti, "HMAC:
  Keyed-Hashing for Message Authentication," RFC 2104, IETF, Feb. 1997.
\item
  \textbf{{[}FIPS198-1{]}} NIST, "The Keyed-Hash Message Authentication
  Code (HMAC)," FIPS PUB 198-1, 2008.
\item
  \textbf{{[}BCK96{]}} M. Bellare, R. Canetti, H. Krawczyk, "Keying Hash
  Functions for Message Authentication," CRYPTO 1996. (HMAC/NMAC PRF
  security; see also M. Bellare, "New Proofs for NMAC and HMAC," CRYPTO
  2006.)
\item
  \textbf{{[}BBT16{]}} M. Bellare, D. J. Bernstein, S. Tessaro,
  "Hash-Function Based PRFs: AMAC and Its Multi-User Security,"
  EUROCRYPT 2016. (Tight multi-user security of HMAC-style PRFs.)
\item
  \textbf{{[}Stripe-sig{]}} Stripe, "Verify webhook signatures,"
  docs.stripe.com/webhooks/signature (accessed 2026-07-04).
\end{itemize}

\begin{quote}
(BCK96/Bellare06), the tight multi-user HMAC bound (BBT16), and the
Stripe timestamped scheme are primary-sourced (verified-facts Part 3;
fetched from primaries, flagged for a final quote-level re-verification
before submission). The tenant-binding and rotation theorems (IV-C/D)
and the replay bound (IV-E) are original results; their proofs reduce
solely to the standard EUF-CMA/PRF assumptions and assert no unverified
external facts.
\end{quote}

\section{V. Lightweight NLP Heuristics for Webhook
Screening}\label{v-lightweight-nlp-heuristics-for-webhook-screening}

Having authenticated the webhook (§IV), the firewall screens its content
for injection signatures before forwarding. This section defends the
\emph{architectural} choice of a deterministic screening layer over an
LLM-based guardrail at this position (§V-A), specifies the Aho--Corasick
signature scanner (§V-B) and the linear feature passes (§V-C), gives
input sanitization (§V-E), accounts the total per-request cost against
the §III budget (§V-F), and states honestly the accuracy ceiling of
heuristics relative to ML classifiers (§V-D).

\subsection{V-A. Why Deterministic Screening at the Edge, Not an LLM
Guardrail}\label{v-a-why-deterministic-screening-at-the-edge-not-an-llm-guardrail}

A natural objection is: \emph{why not run an LLM guardrail (e.g., Llama
Guard, a fine-tuned DeBERTa injection classifier) at the edge instead of
hand-maintained signatures?} The answer is not that guardrails are less
accurate --- for many inputs they are more accurate (§V-D) --- but that
a model-based guardrail is \textbf{structurally incompatible with this
architectural layer}, for four independent reasons rooted in §III.

\begin{enumerate}
\def\labelenumi{\arabic{enumi}.}
\tightlist
\item
  \textbf{Compute placement.} Transformer inference is not a
  sub-millisecond, 128 MB-isolate workload. Llama Guard ($\approx$7--8 B
  parameters) requires GPU-class accelerators and gigabytes of weights;
  even a small encoder-only classifier (DeBERTa-base, $\approx$184 M params)
  exceeds the isolate\textquotesingle s 128 MB memory ceiling (§III-B)
  once weights, activations, and the runtime are counted, and its
  matrix-multiply inference is orders of magnitude beyond a
  \(\kappa N_{\max}\) linear scan. The edge isolate is the \emph{wrong
  hardware} for tensor math.
\item
  \textbf{The guardrail becomes a network hop, not local compute.}
  Deploying the model behind an inference API turns screening into a
  \emph{subrequest} --- adding a full RTT (tens to hundreds of ms of
  wall-clock, per the layered-detector measurements: model deep-scan
  adds $\approx$300--800 ms {[}Layered{]}) to \emph{every} webhook, including
  the benign majority. This inverts the design goal of §I-B: the cheap
  perimeter check now costs more than the origin call it is meant to
  gate. A deterministic scan keeps screening as \emph{local CPU}
  (I/O-free), preserving the wall-clock SLO.
\item
  \textbf{Determinism and auditability.} A signature match is
  explainable and reproducible: the threat log (§VI) records
  \emph{which} signature fired. A model score is neither, complicating
  the non-repudiation obligation (Table I, row 5) and incident
  forensics.
\item
  \textbf{Attack surface.} An LLM guardrail is itself susceptible to
  prompt injection --- the very threat it screens for
  {[}arXiv-2601.07185 finds SFT defenses learn brittle surface
  heuristics{]}. A finite-state automaton has no instruction-following
  semantics to subvert.
\end{enumerate}

The correct architecture is therefore \textbf{layered}: a deterministic,
high-precision, sub-millisecond scan at the edge (this section) that
rejects known-signature attacks cheaply and, for the residual uncertain
traffic, \emph{optionally} escalates to a model-based deep scan at the
origin --- where GPU compute and higher latency budgets exist. This
paper\textquotesingle s contribution is the edge layer; the escalation
tier is discussed as future work (§IX). The edge layer\textquotesingle s
role is not to be the last word on subtle semantic attacks but to
eliminate the high-volume, signature-expressible ones within the CPU
envelope, so that the expensive tiers see less traffic.

\subsection{V-B. Layer 1: Aho--Corasick Multi-Signature
Scan}\label{v-b-layer-1-ahocorasick-multi-signature-scan}

\textbf{Problem.} Given a signature set \(S=\{s_1,\dots,s_p\}\)
(injection phrases, jailbreak markers, control tokens; Appendix C),
determine which \(s_j\) occur in the sanitized payload \(x\) of length
\(n\).

\textbf{Why not per-pattern matching.} Running a separate matcher per
signature costs \(O\!\big(\sum_j (n+|s_j|)\big)=O(pn+m)\) where
\(m=\sum_j|s_j|\) --- linear in the \emph{number} of signatures \(p\).
With \(p\) in the hundreds to thousands, the \(pn\) term dominates and
scales poorly. A single regular expression alternation
\passthrough{\lstinline!(s\_1|...|s\_p)!} compiled to an NFA risks
catastrophic backtracking on adversarial input (a ReDoS vector), which
is itself a DoS surface (Table I, row 7).

\textbf{Aho--Corasick.} The Aho--Corasick automaton is a trie of all
signatures augmented with failure (suffix) links, forming a
deterministic finite automaton that matches \textbf{all} signatures in a
\textbf{single pass} over \(x\), regardless of \(p\) {[}Springer-AC{]}:

\begin{itemize}
\tightlist
\item
  \textbf{Construction:} \(O(m)\) time and space in the total signature
  length \(m\) --- built \textbf{once at module top level} (§III-A) and
  amortized to zero on the request path.
\item
  \textbf{Search:} \(O(n+z)\), where \(n=|x|\) and \(z\) is the number
  of matches reported; each input byte triggers a constant number of
  automaton transitions (goto/failure), with \textbf{no backtracking}.
\end{itemize}

Contrasted with KMP --- a \emph{single}-pattern algorithm at
\(O(n+|s_j|)\) per pattern, hence \(O(pn+m)\) for \(p\) patterns ---
Aho--Corasick removes the dependence on \(p\) from the search cost
entirely {[}Springer-AC{]}. This is the property that makes hundreds of
signatures affordable in one scan.

\textbf{Match cap (bounding \(z\)).} Since \(z\le n\) in the worst case
(e.g., many overlapping short signatures), we cap reported matches at a
constant \(z_{\max}\) and early-exit on the first \emph{blocking}-class
signature, so the effective per-request search cost is
\(O(n+\min(z,z_{\max}))=
O(n)\). Early exit is safe for a \emph{reject} decision (one confirmed
malicious signature suffices) and the cap only limits \emph{enumeration}
for logging, not detection.

\subsection{V-C. Layer 2: Linear Structural and Entropy
Features}\label{v-c-layer-2-linear-structural-and-entropy-features}

Signatures catch \emph{known} phrasings; a small set of \(O(n)\)
structural features catches \emph{structural} injection independent of
exact wording. Each is a single linear pass, contributing to the
constant \(k\le 4\) of §III-C:

\begin{itemize}
\tightlist
\item
  \textbf{Delimiter/role-token density.} Count occurrences of
  instruction-boundary and role markers
  (\passthrough{\lstinline!\#\#\# instruction!},
  \passthrough{\lstinline!<|system|>!},
  \passthrough{\lstinline!assistant:!}, fenced-block openers) normalized
  by length; a spike signals delimiter-spoofing / tag-injection.
  \(O(n)\), computed as a by-product of the Aho--Corasick pass over a
  token sub-dictionary.
\item
  \textbf{Imperative-override heuristic.} Presence of override
  collocations ("ignore/disregard \ldots{} previous/above \ldots{}
  instructions/prompt") captured as multi-word signatures in Layer 1;
  scored as a weighted feature rather than an immediate block to limit
  false positives on benign text that quotes such phrases.
\item
  \textbf{Shannon entropy.}
  \(\hat{H}(x)=-\sum_{c} \hat p(c)\log_2 \hat p(c)\) over the
  byte/character distribution, one pass to accumulate counts, \(O(n)\).
  Anomalously high entropy flags encoded or obfuscated payloads
  (base64/hex smuggling) that evade literal signatures; anomalously low
  entropy with high delimiter density flags templated injection. Entropy
  is a \emph{weak} feature used only to \emph{escalate} (route to the
  origin deep-scan tier), never to block alone (§V-D).
\end{itemize}

The layer emits a score vector; a linear threshold rule combines Layer 1
blocking hits (hard reject), weighted Layer 1/2 features (soft score),
and an escalate/forward/reject decision. All features are \(O(n)\) and
share the input scan, so Layer 2 adds a small constant multiple of
\(n\), not a new asymptotic term.

\subsection{V-D. Accuracy Ceiling: Heuristics vs. ML Classifiers (Honest
Positioning)}\label{v-d-accuracy-ceiling-heuristics-vs-ml-classifiers-honest-positioning}

We do \textbf{not} claim deterministic heuristics match ML classifiers
on subtle or novel injections. The evidence base is explicit about the
ceiling, and we report it rather than obscure it:

\begin{itemize}
\tightlist
\item
  Lexical/deterministic filters achieve moderate standalone
  discrimination --- reported ROC-AUC in the $\approx$0.65 band for
  Aho--Corasick / regex-denylist style detectors on injection corpora
  {[}arXiv-2601.07185{]} \emph{(from source; to be re-verified at claim
  level, see caveat)}.
\item
  Embedding-based ML classifiers can do better: a Random Forest over
  \passthrough{\lstinline!text-embedding-3-small!} features reached AUC
  0.764 / P 0.867 / R 0.870, outperforming encoder-only baselines
  (DeBERTa variants at AUC $\approx$0.50--0.59) {[}arXiv-2410.22284{]} ---
  though this rests on a single preprint with an author-curated set and
  modest absolute AUC (medium confidence, verified-facts Part 2).
\item
  Fine-tuned guardrails achieve high attack-rejection but can learn
  brittle surface correlations with large OOD degradation
  {[}arXiv-2601.07185{]} (medium confidence, single source).
\end{itemize}

The design implication is precisely the layering of §V-A: the edge
heuristics are positioned as a \textbf{high-precision first filter}
(favoring low false-positive rate at a chosen operating point, so benign
traffic is not wrongly blocked) that cheaply removes
signature-expressible attacks, with the \emph{recall gap} on
novel/semantic attacks delegated to an optional origin-side model tier.
§VII measures the edge layer\textquotesingle s own FPR/FNR at its
operating point; we make no claim that the edge layer alone achieves the
full-system detection rate, only that it does so for the
signature-expressible class it targets, within the CPU budget.

\subsection{V-E. Input Sanitization and
Canonicalization}\label{v-e-input-sanitization-and-canonicalization}

Signatures match a \emph{canonical} representation; an attacker who can
vary encoding while preserving meaning evades a naive scan. Before Layer
1, the firewall (Fig. 1, step 3) applies, in one or two linear passes:
(a) \textbf{Unicode normalization} (NFKC) to fold
compatibility/homoglyph variants; (b) \textbf{decoding of transport
encodings} actually in scope (percent-encoding, HTML entities) where the
downstream will decode them, to prevent encode-smuggling past the
scanner; (c) \textbf{whitespace and zero-width-character collapse}
(zero-width joiners/spaces are a known signature-splitting trick); and
(d) \textbf{case folding} for case-insensitive signatures. Sanitization
is bounded by \(O(n)\) and its output length is \(\le c\cdot n\) for a
small constant, preserving the \(N_{\max}\) bound. Crucially,
sanitization is applied to the copy that is \emph{scanned}; the
\emph{forwarded} payload is the original (the signature over which was
already verified in §IV), so normalization cannot alter authenticated
content --- it only informs the block/forward decision.

\subsection{V-F. Per-Request Cost Accounting Against the §III
Budget}\label{v-f-per-request-cost-accounting-against-the-iii-budget}

Collecting the layers, the per-request CPU cost is
\fiteq{C_{\text{req}} \le \underbrace{c_{\text{san}}n}_{\text{§V-E}} +
\underbrace{c_{\text{ac}}(n+z_{\max})}_{\text{§V-B, capped}} +
\underbrace{k\,c_{\text{lin}}n}_{\text{§V-C},\ k\le4} +
\underbrace{c_{\text{hmac}}n}_{\text{§IV}} = \kappa\,n \le \kappa N_{\max},}
exactly the form of §III-C: a constant \(\kappa\) times a bounded input,
with the one-time \(O(m)\) automaton construction excluded (module
scope). There is no super-linear term, no backtracking, and no
per-request allocation proportional to \(p\) (the signature count lives
in the shared, once-built automaton counted against the 128 MB isolate
memory, not per request). The absolute sub-millisecond claim is thus
reduced to a single measured quantity \(\kappa/\rho\) (§III-C),
evaluated at \(N_{\max}=128\) KiB against the 10 ms Free-plan ceiling in
§VII. The structural guarantee --- bounded, strictly linear,
backtrack-free, construction-amortized --- holds regardless of the
measured constant, which is the claim §III promised this section would
substantiate.

\begin{center}\rule{0.5\linewidth}{0.5pt}\end{center}

\subsubsection{Citation keys}\label{citation-keys-3}

\begin{itemize}
\tightlist
\item
  \textbf{{[}Springer-AC{]}} "\emph{(Aho--Corasick multi-pattern
  matching for signature detection)}," Springer,
  doi:10.1007/978-3-031-96093-2\_15.
\item
  \textbf{{[}arXiv-2601.07185{]}} "Defenses Against Prompt Attacks Learn
  Surface Heuristics," arXiv:2601.07185.
\item
  \textbf{{[}arXiv-2410.22284{]}} Ayub \& Majumdar,
  "\emph{(embedding-based prompt-injection classification)}," CAMLIS
  2024, arXiv:2410.22284.
\item
  \textbf{{[}Layered{]}} "\emph{(layered prompt-injection detector;
  \textless1 ms Layers 1--2, 300--800 ms model layer)}," dev.to
  (secondary; used only for the model-tier latency order-of-magnitude).
\end{itemize}

\begin{quote}
heuristic-vs-ML figures (V-D) are MEDIUM confidence / single-source and
are presented as such, with the lexical-filter AUC flagged for
claim-level re-verification. V-A\textquotesingle s LLM-guardrail
incompatibility argument rests on the CONFIRMED §III envelope (128 MB,
sub-ms CPU) plus the model-tier latency order-of-magnitude from
{[}Layered{]} (secondary --- used only qualitatively).
\end{quote}

\section{VI. Asynchronous Threat Logging and Data
Architecture}\label{vi-asynchronous-threat-logging-and-data-architecture}

Every interception decision (Fig. 1) must be recorded to a durable,
queryable, immutable audit trail --- for forensics, for detection
tuning, and for compliance (§VI-F) --- yet the recording must never
appear on the request\textquotesingle s critical path. This section
resolves that tension: §VI-A states the requirements, §VI-B gives the
decoupled write path via \passthrough{\lstinline!waitUntil!} plus a
durable queue, §VI-C--E justify PostgreSQL JSONB with TOAST, WAL/GIN
tuning, and date partitioning for high-velocity ingestion at scale, and
§VI-F maps the design to PCI-DSS and SOC 2 while enforcing data
minimization.

\subsection{VI-A. Requirements}\label{vi-a-requirements}

The threat log must satisfy five properties simultaneously: \textbf{(R1)
non-blocking} --- logging adds zero wall-clock latency to the webhook
response; \textbf{(R2) durable} --- an accepted event is not lost under
isolate eviction or transient DB unavailability; \textbf{(R3)
high-velocity} --- ingestion sustains attack bursts (a flood produces
one log row per rejected request) without back-pressuring the edge;
\textbf{(R4) queryable} --- heterogeneous, schema-varying attack
payloads remain indexable for analysis; and \textbf{(R5) immutable and
compliant} --- the record is append-only and admissible as an audit
trail without itself becoming a data-exposure liability. R1 and R2 are
in tension (durability usually implies waiting); §VI-B resolves it by
moving durability off the request path into a queue.

\subsection{\texorpdfstring{VI-B. Decoupled Execution:
\texttt{waitUntil} + Durable
Queue}{VI-B. Decoupled Execution: waitUntil + Durable Queue}}\label{vi-b-decoupled-execution-waituntil--durable-queue}

\textbf{\passthrough{\lstinline!waitUntil!} mechanics.} The Workers
runtime exposes \passthrough{\lstinline!ctx.waitUntil(promise)!} (the
successor to \passthrough{\lstinline!FetchEvent.waitUntil!}), which
registers a promise whose completion the runtime awaits \emph{after the
\passthrough{\lstinline!Response!} has already been returned to the
client}. Concretely, the handler computes the block/forward decision
(Fig. 1, steps 1--6), returns the response (step forward/reject), and
calls \passthrough{\lstinline!ctx.waitUntil(logThreat(event))!}; the
isolate is kept alive to drain \passthrough{\lstinline!logThreat!} but
the client\textquotesingle s latency is bounded by the response, not by
the log write. Because the log write is network I/O, it \textbf{does not
consume the CPU budget} (§III-B) --- R1 is met by construction, and the
sub-millisecond CPU analysis of §V-F is unaffected by logging.

\textbf{Why \passthrough{\lstinline!waitUntil!} alone is insufficient
for R2.} \passthrough{\lstinline!waitUntil!} is \emph{best-effort}: if
the isolate is evicted or the write fails, the event is lost, and it
offers no back-pressure or batching. Writing directly to PostgreSQL from
the edge also couples every request to a database connection ---
untenable at edge concurrency and a head-of-line risk if the DB slows.
We therefore interpose a \textbf{durable queue} (Cloudflare Queues, or a
Durable Object buffer): \passthrough{\lstinline!waitUntil!} performs a
single fast \passthrough{\lstinline!queue.send(event)!} (one
subrequest), and a separate \textbf{queue consumer} Worker batches
events and performs the PostgreSQL insert. This yields:

\begin{itemize}
\tightlist
\item
  \textbf{R1/R2 both:} the edge pays only a queue enqueue (fast,
  durable-once-acked); durability and retry live in the consumer, off
  the request path.
\item
  \textbf{Batching for R3:} the consumer inserts in batches (multi-row
  \passthrough{\lstinline!INSERT!} / \passthrough{\lstinline!COPY!}),
  amortizing per-row WAL and round-trip cost --- the single most
  effective throughput lever (§VI-D).
\item
  \textbf{Back-pressure isolation:} a DB slowdown grows the queue, not
  the webhook latency; the edge never blocks. The queue depth is itself
  a monitorable DoS signal.
\end{itemize}

Delivery is at-least-once (consumer retries on failure), so the schema
uses an idempotency key (the authenticated nonce \(\eta\) from §IV, or a
synthetic event id) with
\passthrough{\lstinline!INSERT ... ON CONFLICT DO NOTHING!} to dedupe
replays of the log write.

\subsection{VI-C. JSONB for Unstructured, High-Velocity Attack
Payloads}\label{vi-c-jsonb-for-unstructured-high-velocity-attack-payloads}

Attack payloads are heterogeneous: different injection classes carry
different fields, and the raw webhook body has no fixed schema across
tenants. A rigid relational schema would require migration per new
attack shape; a text blob is not queryable (R4). PostgreSQL
\passthrough{\lstinline!jsonb!} resolves this:

\begin{itemize}
\tightlist
\item
  \textbf{Binary, decomposed storage.} Per the official documentation,
  \passthrough{\lstinline!json!} stores an exact text copy that
  "processing functions must reparse on each execution," whereas
  \passthrough{\lstinline!jsonb!} is stored "in a decomposed binary
  format that makes it slightly slower to input due to added conversion
  overhead, but significantly faster to process, since no reparsing is
  needed" {[}PG-json{]}. For a write-once/read-many audit log queried
  during investigations, the one-time input cost is worth the repeated
  query speedup, and --- decisively for R4 ---
  \textbf{\passthrough{\lstinline!jsonb!} supports indexing} (GIN,
  §VI-D) whereas \passthrough{\lstinline!json!} does not.
\item
  \textbf{TOAST for oversized payloads.} A large injection payload
  (e.g., a multi-kilobyte obfuscated prompt) would otherwise threaten
  the 8 KB heap page. TOAST --- "The Oversized-Attribute Storage
  Technique" --- triggers automatically "only when a row value to be
  stored in a table is wider than
  \passthrough{\lstinline!TOAST\_TUPLE\_THRESHOLD!} bytes (normally 2
  kB)," compressing and/or moving field values out-of-line "until the
  row value is shorter than
  \passthrough{\lstinline!TOAST\_TUPLE\_TARGET!} bytes (also normally 2
  kB)" {[}PG-toast{]}. The \passthrough{\lstinline!jsonb!} column uses
  the default \passthrough{\lstinline!EXTENDED!} strategy (compression
  then out-of-line) {[}PG-toast{]}, so wide payloads are transparently
  compressed and relocated to the TOAST relation, keeping the main-heap
  tuples narrow and the table scannable. This is why unstructured,
  occasionally-large payloads do \textbf{not} bloat the primary heap:
  oversized values live out-of-line by construction. \emph{(Precision
  note: "2 kB" is the documented nominal threshold; we state it as $\approx$2
  kB, not an exact 2048, per the docs\textquotesingle{} "normally 2 kB"
  wording.)}
\end{itemize}

\subsection{VI-D. WAL, Insert Throughput, and Avoiding GIN Index
Bloat}\label{vi-d-wal-insert-throughput-and-avoiding-gin-index-bloat}

\textbf{WAL and the durability/throughput trade.} Every insert is first
written to the Write-Ahead Log for durability. Two knobs govern the
cost, both documentation-confirmed:

\begin{itemize}
\tightlist
\item
  \textbf{\passthrough{\lstinline!synchronous\_commit!}.} Default
  \passthrough{\lstinline!on!} makes each commit wait for WAL flush to
  disk. Setting it \passthrough{\lstinline!off!} removes the wait; per
  the docs this risks losing recent committed transactions (maximum
  delay three times \passthrough{\lstinline!wal\_writer\_delay!}) but
  "does not create any risk of database inconsistency"
  {[}PG-wal-conf{]}. For a \emph{threat log}, this trade is appropriate:
  losing the last few hundred milliseconds of log rows under a crash is
  acceptable (the queue\textquotesingle s at-least-once retry recovers
  most), and there is no cross-row invariant to violate. We set
  \passthrough{\lstinline!synchronous\_commit = off!} for the log
  database (not for any transactional tenant data).
\item
  \textbf{Group commit.} \passthrough{\lstinline!commit\_delay!}
  (default 0) "adds a time delay before a WAL flush \ldots{} allowing a
  larger number of transactions to commit via a single WAL flush,"
  active when at least \passthrough{\lstinline!commit\_siblings!}
  (default 5) transactions are concurrent {[}PG-wal-conf{]}. With
  batched inserts from the consumer this further amortizes flush cost.
\item
  \textbf{Checkpoints.} Checkpoints occur every
  \passthrough{\lstinline!checkpoint\_timeout!} (default 5 min) or when
  \passthrough{\lstinline!max\_wal\_size!} (default 1 GB) is about to be
  exceeded {[}PG-wal-config{]}. Under sustained insert load we raise
  \passthrough{\lstinline!max\_wal\_size!} and keep
  \passthrough{\lstinline!checkpoint\_completion\_target!} at its
  default 0.9 (the recommended maximum) to spread checkpoint I/O and
  avoid write stalls {[}PG-wal-config{]}. With
  \passthrough{\lstinline!full\_page\_writes!} on (torn-page
  protection), frequent checkpoints inflate WAL volume, so fewer,
  wider-spaced checkpoints favor insert throughput.
\end{itemize}

\textbf{The GIN index-bloat problem and its mitigation.} To keep
payloads queryable (R4) we index the \passthrough{\lstinline!jsonb!}
column with GIN. But GIN updates are "inherently slow \ldots{} inserting
or updating one heap row can cause many inserts into the index (one for
each key extracted from the indexed item)" {[}PG-gin{]} --- exactly the
pathology that would throttle high-velocity ingestion.
PostgreSQL\textquotesingle s \passthrough{\lstinline!fastupdate!}
mechanism defers this: new entries go "into a temporary, unsorted list
of pending entries," flushed to the main GIN structure on
vacuum/autoanalyze, on
\passthrough{\lstinline!gin\_clean\_pending\_list()!}, or when the list
exceeds \passthrough{\lstinline!gin\_pending\_list\_limit!}
{[}PG-gin{]}. We therefore (i) enable
\passthrough{\lstinline!fastupdate!} and size
\passthrough{\lstinline!gin\_pending\_list\_limit!} so pending-list
flushes coincide with batch boundaries, turning per-row index
maintenance into periodic bulk maintenance; and (ii) index only the
\passthrough{\lstinline!jsonb!} sub-paths actually queried (e.g., a
\passthrough{\lstinline!jsonb\_path\_ops!} GIN on the signature/verdict
keys) rather than the whole document, reducing extracted-key fan-out.
\emph{(The \passthrough{\lstinline!jsonb\_path\_ops!} vs
\passthrough{\lstinline!jsonb\_ops!} operator-class trade is in scope
but was not primary-verified; see evidence status.)}

\subsection{VI-E. Native Declarative Partitioning for Scale and
Retention}\label{vi-e-native-declarative-partitioning-for-scale-and-retention}

A monotonically growing log table degrades: index size grows, autovacuum
lengthens, and time-window queries scan irrelevant history.
PostgreSQL\textquotesingle s \textbf{native declarative partitioning}
(built in since PG 10, distinct from and more performant than legacy
inheritance/\passthrough{\lstinline!UNION ALL!} {[}PG-partition{]})
addresses all three. We declare
\passthrough{\lstinline!PARTITION BY RANGE (created\_at)!} with one
partition per day (or week), bounds "inclusive at the lower end and
exclusive at the upper end" {[}PG-partition{]}:

\begin{itemize}
\tightlist
\item
  \textbf{Ingestion locality.} Inserts route automatically to the
  current partition {[}PG-partition{]}; its indexes stay small and
  cache-resident, so index maintenance cost is bounded by
  \emph{today\textquotesingle s} volume, not all history --- directly
  countering the bloat of §VI-D at the table level.
\item
  \textbf{Query pruning.} Investigations are time-scoped ("attacks in
  the last 24 h"); partition pruning restricts the scan to the relevant
  partitions.
\item
  \textbf{O(1) retention.} Aging out old logs uses
  \passthrough{\lstinline!DETACH PARTITION!}/\passthrough{\lstinline!DROP TABLE!},
  which the docs note is "far faster than a bulk operation" and
  "entirely avoid{[}s{]} the VACUUM overhead caused by a bulk DELETE"
  {[}PG-partition{]}. Retention rollover (§VI-F) becomes a metadata
  operation, not a billion-row delete.
\item
  \textbf{Constraint.} Any primary/unique key on a partitioned table
  "must include all of the partition key columns" {[}PG-partition{]};
  our idempotency key is therefore
  \passthrough{\lstinline!(created\_at, event\_id)!}, which both
  satisfies the constraint and aligns dedup with the partition.
\end{itemize}

\subsection{VI-F. Compliance Guardrails: PCI-DSS, SOC 2, and Data
Minimization}\label{vi-f-compliance-guardrails-pci-dss-soc-2-and-data-minimization}

An immutable attack-log serves compliance, but a naïvely-implemented one
\emph{becomes} a breach: an attack payload aimed at a payment gateway
may itself contain a Primary Account Number (PAN) or PII. The design
must satisfy the audit-trail mandate \textbf{and} the data-minimization
mandate at once.

\textbf{Audit-trail mapping.} PCI-DSS v4.0.1 \textbf{Requirement 10}
("Log and monitor all access to system components and cardholder data")
governs the audit trail; sub-requirement \textbf{10.5.1} mandates
retaining audit-log history for \textbf{at least 12 months}, with the
\textbf{most recent 3 months immediately available} for analysis ("hot"
storage) {[}PCI-DSS{]}. SOC 2 maps to the AICPA Trust Services Criteria
Common Criteria family \textbf{CC7}: \textbf{CC7.2} (monitor system
components to detect anomalies and security events) and \textbf{CC7.3}
(evaluate security events and trigger incident response), CC7.1
(vulnerability detection), and CC7.4 (incident response)
{[}AICPA-TSC{]}. An immutable, asynchronously-written PostgreSQL audit
trail supplies exactly the \emph{detection and record} evidence these
criteria require; unlike PCI-DSS, SOC 2 fixes no statutory retention
period (it is auditor/period-defined). \textbf{Immutability} is enforced
architecturally, satisfying PCI-DSS \textbf{10.3.2} (audit logs
protected from modification) and \textbf{10.3.1} (read access limited to
a job-related need): the log role has \passthrough{\lstinline!INSERT!}
and \passthrough{\lstinline!SELECT!} privileges only --- no
\passthrough{\lstinline!UPDATE!}/\passthrough{\lstinline!DELETE!} --- so
records are append-only; the \emph{sole} deletion path is time-based
partition \passthrough{\lstinline!DROP!} (§VI-E), executed by a separate
retention role on a fixed schedule, which is precisely the controlled,
policy-driven expiry the retention clause expects rather than ad hoc row
deletion. What to log is governed by \textbf{10.2} (logs sufficient to
support anomaly detection and forensics); the firewall records every
interception decision to that end.

\textbf{Data minimization (the critical guardrail).} We never persist
raw sensitive data. Before the edge enqueues an event (§VI-B), a bounded
redaction pass (an extension of the §V-E sanitization, \(O(n)\)) masks
high-confidence sensitive tokens: PAN candidates (Luhn-valid digit runs)
are replaced by a truncated token (first-6/last-4 with the middle masked
--- exactly the maximum display exposure PCI-DSS \textbf{3.4.1} permits,
BIN + last four --- or a keyed hash rendering the value unreadable per
\textbf{3.5.1}, which lists one-way hashing, truncation, index tokens,
or strong cryptography), and recognizable PII patterns are masked.
Sensitive authentication data (full track, CVV, PIN) is never retained
at all, per \textbf{3.3}. What the log stores is therefore the
\textbf{attack metadata} --- which signatures fired (§V-B), the feature
scores (§V-C), the verdict, tenant id, timestamps --- and a
\textbf{redacted} payload sufficient for forensic pattern analysis but
stripped of cardholder data. This reconciles "log the attack" with
"store no CHD": the forensic value of an injection payload lies in its
\emph{instruction structure}, which survives redaction, not in any
incidental PAN it carries, which does not need to. Redaction at the
\emph{edge} (before the queue) ensures raw CHD never even reaches the
logging tier, minimizing the systems in PCI scope.

\begin{quote}
\textbf{Evidence status for §VI-F:} the control references --- PCI-DSS
v4.0.1 \textbf{10.5.1} ("Retain audit log history for at least 12
months, with at least the most recent three months immediately available
for analysis," verbatim standard text), \textbf{10.3.1/10.3.2} (log
protection), \textbf{10.2} (logging scope), \textbf{3.3} (no SAD after
authorization), \textbf{3.4.1} (PAN display max BIN + last 4),
\textbf{3.5.1} (stored PAN rendered unreadable); and AICPA TSC
\textbf{CC7.1--CC7.4} --- are all CONFIRMED against the primary PCI SSC
v4.0.1 standard and the AICPA 2017 TSC (2022 revised) in a dedicated
verification pass. §VI-F is now on the same primary-sourced footing as
the rest of §VI.
\end{quote}

\begin{center}\rule{0.5\linewidth}{0.5pt}\end{center}

\subsubsection{Citation keys (PostgreSQL claims: all CONFIRMED 3-0,
verified-facts Part
4)}\label{citation-keys-postgresql-claims-all-confirmed-3-0-verified-facts-part-4}

\begin{itemize}
\tightlist
\item
  \textbf{{[}PG-json{]}} PostgreSQL Global Dev. Group, "JSON Types,"
  postgresql.org/docs/current/datatype-json.html.
\item
  \textbf{{[}PG-toast{]}} "Database Physical Storage --- TOAST,"
  postgresql.org/docs/current/storage-toast.html.
\item
  \textbf{{[}PG-wal-conf{]}} "Write Ahead Log --- runtime config,"
  postgresql.org/docs/current/runtime-config-wal.html.
\item
  \textbf{{[}PG-wal-config{]}} "WAL Configuration,"
  postgresql.org/docs/current/wal-configuration.html.
\item
  \textbf{{[}PG-gin{]}} "GIN Indexes,"
  postgresql.org/docs/current/gin.html.
\item
  \textbf{{[}PG-partition{]}} "Table Partitioning,"
  postgresql.org/docs/current/ddl-partitioning.html.
\item
  \textbf{{[}PCI-DSS{]}} PCI Security Standards Council, "Payment Card
  Industry Data Security Standard v4.0.1," June 2024. Req 10.2
  (logging), 10.3.1/10.3.2 (log protection), 10.5.1 (12-month retention
  / 3-month hot); Req 3.3 (no SAD), 3.4.1 (PAN display masking), 3.5.1
  (PAN rendered unreadable).
\item
  \textbf{{[}AICPA-TSC{]}} AICPA, "Trust Services Criteria for Security,
  Availability, Processing Integrity, Confidentiality, and Privacy"
  (2017, rev. 2022), Common Criteria CC7.1--CC7.4.
\end{itemize}

\begin{quote}
\passthrough{\lstinline!synchronous\_commit!}/\passthrough{\lstinline!commit\_delay!}/checkpoint/\passthrough{\lstinline!full\_page\_writes!}
behavior, GIN \passthrough{\lstinline!fastupdate!}/
\passthrough{\lstinline!gin\_pending\_list\_limit!}, and native RANGE
partitioning with DETACH/DROP retention --- are all CONFIRMED at high
confidence (verified-facts Part 4, 25/25 claims 3-0, official docs).
\passthrough{\lstinline!waitUntil!} semantics (§VI-B) follow the Workers
execution model of §III (CONFIRMED). §VI-F compliance clauses (PCI Req
10.2/10.3.1/10.3.2/10.5.1, 3.3/3.4.1/3.5.1; SOC 2 CC7.1--CC7.4) are
CONFIRMED against the primary PCI SSC v4.0.1 standard and AICPA 2017
TSC. The \passthrough{\lstinline!jsonb\_path\_ops!} choice (§VI-D) is
noted as not-yet-verified.
\end{quote}

\section{VII. Empirical Evaluation}\label{vii-empirical-evaluation}

This section defines the evaluation of two claims the architecture must
support: \textbf{latency} --- the edge firewall adds bounded, sub-SLO
latency at the tail (§VII-B analytical, §VII-C protocol) --- and
\textbf{detection efficacy} --- the edge layer mitigates the
signature-expressible attack class at high rate with a controlled
false-positive rate (§VII-D). We first fix the methodology and corpus
(§VII-A), then give the analytical latency model, the measurement
protocols, the statistical treatment (§VII-E), and threats to validity
(§VII-F).

\begin{quote}
\textbf{Reporting convention.} This paper specifies the experiment and
does \textbf{not} assert result values it has not measured.
\textbf{Detection-efficacy results (§VII-D) are populated from a real
run} over the 662-row deepset corpus via the Phase-2 harness;
\textbf{latency percentiles (§VII-C) remain placeholders}
\(\langle\cdot\rangle\) pending a full edge deployment. Analytical
bounds (§VII-B) are derived and labelled as such, distinct from
empirical percentiles. Where a run overturned an initial expectation ---
as the detection numbers do for the "\(>99\%\)" ambition --- we report
the measurement and revise the claim, not the reverse.
\end{quote}

\subsection{VII-A. Methodology and Benchmark
Corpus}\label{vii-a-methodology-and-benchmark-corpus}

\textbf{Testbed.} The firewall Worker is deployed to the edge runtime
(Cloudflare Workers) with a mock origin that returns a fixed-size
response after a controlled service time \(S_O\), isolating the
firewall\textquotesingle s \emph{added} latency from origin variability.
Load is generated with \textbf{k6} (distributed, scriptable, native
percentile output) cross-checked against \textbf{autocannon} for a
second measurement path. Each configuration is run for \(R\) repetitions
after a warm-up that guarantees a warm isolate (§III-A), so cold-start
is measured separately (§VII-C) rather than contaminating the
steady-state distribution.

\textbf{Benchmark corpus (documented provenance).} Rather than the
folklore "200-payload set" --- which corresponds to no real dataset ---
we use a corpus of \emph{verified provenance}:

\begin{itemize}
\tightlist
\item
  \textbf{Primary:} \passthrough{\lstinline!deepset/prompt-injections!}
  (Hugging Face) --- \textbf{662 examples} (546 train / 116 test),
  binary labels \textbf{0 = legitimate, 1 = injection}, columns
  \passthrough{\lstinline!text!}/\passthrough{\lstinline!label!},
  German+English, Apache 2.0 {[}DS-PI{]}. The \textbf{malicious payload
  set} \(\mathcal{X}^{+}\) is the \passthrough{\lstinline!label==1!}
  subset; the \textbf{benign control set} \(\mathcal{X}^{-}\) (for
  false-positive measurement) is the \passthrough{\lstinline!label==0!}
  subset.
\item
  \textbf{Diversity supplement:}
  \passthrough{\lstinline!jackhhao/jailbreak-classification!} (1,306
  rows,
  \passthrough{\lstinline!jailbreak!}/\passthrough{\lstinline!benign!},
  Apache 2.0 {[}DS-JB{]}) and
  \passthrough{\lstinline!Lakera/gandalf\_ignore\_instructions!}
  {[}DS-GA{]}, to test recall on jailbreak and instruction-ignore
  families beyond the primary set.
\end{itemize}

The exact counts \(|\mathcal{X}^{+}|, |\mathcal{X}^{-}|\) used in each
experiment are reported \textbf{from the data at run time} (a
seed-fixed, documented split), not assumed; the harness prints them for
inclusion in the results table. This makes the corpus reproducible from
named, licensed sources --- directly remedying the unsupported dataset
provenance the design review flagged.

\subsection{VII-B. Analytical Latency Model (the theoretical
bound)}\label{vii-b-analytical-latency-model-the-theoretical-bound}

We decompose the firewall\textquotesingle s \emph{added} end-to-end
latency (excluding origin service time \(S_O\), which the firewall does
not control) into its constituent terms, per the Fig. 1 pipeline:
\fiteq{L_{\text{added}} \;=\; \underbrace{L_{\text{CPU}}^{\text{proc}}}_{\text{steps 1--6, local}}
\;+\; \underbrace{L_{\text{enq}}}_{\text{waitUntil enqueue, off-path}}
\;+\; \underbrace{L_{\text{net}}^{\Delta}}_{\text{proxy hop overhead}}.}

\textbf{The CPU term.} From §III-C/§V-F, the local processing cost is
\(C_{\text{req}}\le\kappa N_{\max}\) operations, i.e. wall-time
\(L_{\text{CPU}}^{\text{proc}} = C_{\text{req}}/\rho \le
(\kappa/\rho)\,N_{\max}\). This is the term the theory bounds;
\(\kappa/\rho\) (per-byte wall cost) is the single empirical constant,
measured in §VII-C. Structurally it is \(O(N_{\max})\) with no
super-linear component, so it cannot exhibit a heavy tail from the
algorithm itself.

\textbf{The off-path term.} \(L_{\text{enq}}\) is the
\passthrough{\lstinline!queue.send!} enqueue (§VI-B). Because the
threat-log write is deferred via \passthrough{\lstinline!waitUntil!}
\emph{after} the response returns, \(L_{\text{enq}}\) contributes to the
\emph{response} path only as a single fast subrequest and the DB write
contributes \textbf{zero} (it happens after the client is served). We
include \(L_{\text{enq}}\) conservatively; if the enqueue itself is also
deferred behind the response, this term vanishes from
\(L_{\text{added}}\).

\textbf{The proxy-hop term.} \(L_{\text{net}}^{\Delta}\) is the
incremental network cost of interposing the edge proxy versus a direct
origin call. Since the edge runtime terminates the connection at a
point-of-presence geographically near the client,
\(L_{\text{net}}^{\Delta}\) is typically \emph{small or negative} (edge
TLS termination can reduce client-perceived latency), but we treat it as
a non-negative unknown bounded by measurement.

\textbf{Tail (p99) argument.} The p99 of \(L_{\text{added}}\) is bounded
by the p99 of each additive term (subadditivity of quantiles does not
hold in general, so we bound conservatively by the sum of per-term p99s
under independence, and validate empirically):
\fiteq{L_{\text{added}}^{p99} \;\lesssim\; \frac{\kappa}{\rho}N_{\max}
\;+\; L_{\text{enq}}^{p99} \;+\; L_{\text{net}}^{\Delta,p99}.} The
\textbf{claim to be tested} is
\(L_{\text{added}}^{p99} < 50\text{ ms}\). The analysis shows the
\emph{algorithmic} term is a bounded constant (sub-millisecond for
realistic \(\kappa/\rho\) and \(N_{\max}=128\) KiB, since even \(\rho\)
on the order of \(10^{8}\) byte-ops/s gives
\(\frac{\kappa}{\rho}N_{\max}\ll 1\) ms), so any violation of the 50 ms
target must originate in the \emph{network/queue} terms, not the
firewall\textquotesingle s computation --- which is the
decomposition\textquotesingle s diagnostic value. §VII-C tests the bound
and attributes any tail to the correct term.

\subsection{VII-C. Empirical Latency
Protocol}\label{vii-c-empirical-latency-protocol}

\begin{itemize}
\tightlist
\item
  \textbf{Percentiles.} Report
  \(\langle p50\rangle, \langle p90\rangle, \langle p99\rangle\) (and
  \(\langle p99.9\rangle\)) of \(L_{\text{added}}\), computed from
  \(\ge N_{\text{req}}\) requests per configuration with the origin
  mock\textquotesingle s \(S_O\) subtracted. Percentiles are estimated
  with the harness\textquotesingle s streaming estimator and
  cross-checked by a bootstrap over raw samples (§VII-E).
\item
  \textbf{Isolating \(\kappa/\rho\).} Sweep payload size
  \(N\in\{1\text{ KiB},\dots,N_{\max}\}\) and regress
  \(L_{\text{CPU}}^{\text{proc}}\) on \(N\); the slope is
  \(\kappa/\rho\) (closing the §III-C/§V-F loop with a measured
  constant), the intercept is fixed overhead. Report
  \(\langle\kappa/\rho\rangle\) with a confidence band.
\item
  \textbf{Cold vs warm.} Report cold-start added latency
  \(\langle L_{\text{cold}}\rangle\) separately (first request to a
  fresh isolate, including module-scope automaton build, §III-A) from
  warm steady state, since they have different distributions and the
  automaton build is a one-time cost.
\item
  \textbf{Load levels.} Repeat at increasing arrival rates to the
  saturation point; report the rate at which \(L_{\text{added}}^{p99}\)
  crosses the SLO, characterizing headroom.
\end{itemize}

\subsection{VII-D. Detection-Efficacy
Methodology}\label{vii-d-detection-efficacy-methodology}

\textbf{Confusion matrix.} Running
\(\mathcal{X}^{+}\cup\mathcal{X}^{-}\) through the firewall at a fixed
decision threshold \(\theta\) (§V-C) yields counts \(TP, FP, TN, FN\).
We report:
\fiteq{\text{FNR}=\frac{FN}{TP+FN},\quad \text{FPR}=\frac{FP}{FP+TN},\quad
\text{Mitigation}=\text{Recall}=\frac{TP}{TP+FN}=1-\text{FNR},\quad
\text{Precision}=\frac{TP}{TP+FP}.} The "\(>99\%\) mitigation" target
is thus the hypothesis \(\text{Recall}>0.99\) \textbf{on the
signature-expressible malicious class} (the scoping of §V-D --- we do
not claim it for arbitrary semantic attacks). It is reported as
\(\langle\text{Recall}\rangle\) with a confidence interval, not
asserted.

\textbf{Interval estimation.} Because \(TP,FP\) are binomial counts,
point rates are insufficient; we report \textbf{Wilson score 95\%
confidence intervals} for each proportion (Wilson is preferred over
normal-approximation for proportions near 0 or 1, exactly the regime of
a \(>99\%\) recall claim). A claim "\(\text{Recall}>0.99\)" is credited
only if the \emph{lower} Wilson bound exceeds 0.99, which in turn
dictates the minimum \(|\mathcal{X}^{+}|\) needed (§VII-E).

\textbf{Operating point and ROC.} We sweep \(\theta\) to trace the
ROC/precision-recall curve and select the operating point by the design
priority of §V-D (high precision / low FPR first filter), and report the
chosen \(\theta\), its FPR/FNR, and the AUC \(\langle\text{AUC}\rangle\)
for comparability with the literature band (lexical filters $\approx$0.65;
embedding classifiers $\approx$0.76, §V-D).

\textbf{Per-family breakdown.} Report recall separately per attack
family (direct override, role-play, delimiter/tag, control-token,
§II-C/Appendix C) and per source dataset, since an aggregate rate can
mask a blind spot in one family --- a per-family table is more
informative and more honest than a single headline number.

\textbf{Results (measured, full deepset corpus).} Running the complete
662-row corpus (\(|\mathcal{X}^{+}|=263\) injection,
\(|\mathcal{X}^{-}|=399\) legitimate) through the shipped edge screening
pipeline at \(\theta=1.0\) yields Table III.

\textbf{TABLE III. Edge screening layer on deepset/prompt-injections
(Wilson 95\% CI).}

{\def\LTcaptype{none} % do not increment counter
\begin{table*}[!t]\centering\scriptsize
\resizebox{\textwidth}{!}{%
\begin{tabular}{@{}lll@{}}
\toprule
Metric & Value & 95\% CI \\
\midrule



Mitigation (Recall) & \textbf{2.28 \%} & {[}1.05 \%, 4.89 \%{]} \\
False-Positive Rate & \textbf{0.00 \%} & {[}0.00 \%, 0.95 \%{]} \\
Precision & 100.00 \% & {[}60.97 \%, 100 \%{]} \\
ROC AUC & \textbf{0.511} & --- \\
\bottomrule
\end{tabular}}
\end{table*}
}

This result is decisive and, we argue, \emph{strengthens} the paper
rather than undermining it, because it confirms the §V-D positioning
empirically and refutes an overclaim. Three readings follow.
\textbf{(i)} The edge layer is a \textbf{high-precision, low-recall
first filter}: it blocked \(6/263\) injections while raising
\textbf{zero} false positives over \(399\) legitimate payloads --- it
never harms benign traffic, the property a perimeter filter must
guarantee. \textbf{(ii)} Its recall on the \emph{full}, diverse,
partly-German deepset distribution is low (\(2.28\%\), AUC \(\approx\)
chance) because a compact deterministic signature set cannot, even in
principle, cover novel or multilingual phrasings --- precisely the
ceiling §V-D anticipated and the motivation for the origin escalation
tier (§IX). \textbf{(iii)} Consequently, \textbf{the "\(>99\%\)
mitigation" figure is a \emph{full-system} target, not an edge-layer
property}, and the data show the signature layer alone does not meet it:
the Wilson lower bound is \(1.05\%\), nowhere near \(0.99\), so the
hypothesis \(\text{Recall}>0.99\) is \textbf{rejected for the edge layer
in isolation}. We therefore restate the system\textquotesingle s
detection claim precisely: the edge layer contributes
\emph{high-precision, zero-false-positive} removal of the
signature-expressible subclass within the sub-millisecond CPU budget,
and the aggregate mitigation rate is a property of the \emph{composed}
edge + origin-tier system, whose measurement we scope to future work
(§IX). Reporting Table III rather than curating a benchmark to flatter
the signature set is the methodological commitment of
§VII\textquotesingle s opening convention.

\textbf{Full-system result (edge + origin escalation tier).} We
implemented and \emph{deployed} the origin-side escalation tier (§IX) as
a self-hosted injection classifier, and measured the \emph{composed}
system on the same corpus --- a payload is blocked iff the edge
hard-rejects it \textbf{or} the classifier flags it. We report two
classifiers run on the live deployment (CPU inference): the ungated
\passthrough{\lstinline!protectai/deberta-v3-base-prompt-injection-v2!}
(English-only), and the chosen, purpose-built multilingual \textbf{Meta
Llama Prompt Guard 2 86M} (Table IV).

\textbf{TABLE IV. Full-system detection on deepset, live deployment
(Wilson 95\% CI).}

{\def\LTcaptype{none} % do not increment counter
\begin{table*}[!t]\centering\scriptsize
\resizebox{\textwidth}{!}{%
\begin{tabular}{@{}llll@{}}
\toprule
Configuration & Recall (Mitigation) & FPR & Precision \\
\midrule



Edge only (Table III) & 2.28 \% {[}1.05, 4.89{]} & 0.00 \% & 100 \% \\
+ protectai (EN-only) & \textbf{41.44 \%} {[}35.66, 47.48{]} & 1.00 \%
{[}0.39, 2.55{]} & 96.46 \% \\
+ Prompt Guard 2 86M (chosen, multilingual) & \textbf{22.81 \%}
{[}18.15, 28.26{]} & \textbf{0.25 \%} {[}0.04, 1.41{]} & 98.36 \% \\
\bottomrule
\end{tabular}}
\end{table*}
}

Four findings follow, and the third \textbf{corrects an expectation
stated in an earlier draft of this work}. \textbf{(i)} The escalation
tier lifts recall an order of magnitude over the edge alone (2.28 \% $\to$
22.8--41.4 \%) --- the layered architecture\textquotesingle s value is
\emph{empirically demonstrated}, not merely argued. \textbf{(ii)} Vendor
accuracy numbers do \textbf{not} transfer across distributions:
protectai self-reports 99.74 \% recall and Prompt Guard 2 self-reports
97.5 \% recall at 1 \% FPR, both \emph{on their own held-out sets}, yet
on this out-of-distribution, partly-German corpus they measure 41.4 \%
and 22.8 \% respectively --- a stark, quantified vindication of
independent evaluation over cited model-card figures. \textbf{(iii)}
Counterintuitively, the flagship multilingual model achieves
\emph{lower} recall than the English-only proxy, because Meta
deliberately tuned Prompt Guard 2 to minimize false positives: it trades
recall for precision, delivering 0.25 \% FPR / 98.4 \% precision versus
the proxy\textquotesingle s 1.0 \% / 96.5 \%. Model selection at this
tier is therefore an \emph{operating-point} decision (precision-first vs
recall-first), not a strict quality ordering --- a nuance invisible to
anyone trusting a single headline accuracy figure. \textbf{(iv)} Even
the chosen model at its default operating point leaves a large residual
recall gap; the "\textgreater99 \% mitigation" ambition is thus a
genuinely \textbf{open problem} (§IX) that dropping in a purpose-built
guardrail does \emph{not} close --- threshold tuning, ensembling (e.g.
adding NVIDIA NemoGuard on the jailbreak axis), and adaptive-attack
hardening remain necessary.

\textbf{Testing those two remedies (threshold sweep + jailbreak-axis
ensemble).} We evaluated both directly. Sweeping Prompt Guard
2\textquotesingle s decision threshold across \([0.9, 0.1]\) moves
recall only from 20.5 \% to \textbf{26.6 \%} (the aggressive end costing
a doubling of FPR to 0.50 \%) --- the recall gap is \textbf{structural,
not a threshold artifact}: PG2 scores \textasciitilde73 \% of deepset
injections \emph{below 0.1}. Adding NVIDIA NemoGuard JailbreakDetect as
an OR ensemble yielded \textbf{0 \% additional recall} on deepset, for
two separable reasons we report honestly: (a) NemoGuard is a
\emph{jailbreak} detector and deepset is \emph{injection} --- an axis
mismatch --- and (b) we could not reproduce NVIDIA\textquotesingle s
exact CPU embedding pipeline from the public artifacts and
self-contradictory model card (a blatant jailbreak scored
\textasciitilde0.09, not \textasciitilde0.99, across every
embedder/pooling/normalization variant tried), so this 0 \% reflects
\emph{our reproduction}, not NemoGuard\textquotesingle s true NIM
performance. The operational conclusion (details in
\passthrough{\lstinline!docs/nemoguard\_ensemble.md!}) is that neither
remedy closes the gap on this corpus; a stronger injection-specific
classifier or a fine-tune on injection data is required.

\textbf{Phase 5 --- the gap was never structural; it is a decision-head
artifact (SOLVED).} We tested whether PG2\textquotesingle s \emph{frozen
encoder} already separates injections, training a lightweight logistic
head on its penultimate (pooler) embeddings --- the exact representation
PG2\textquotesingle s own head classifies --- using a license-clean
corpus (deepset+gandalf+OpenOrca train; §refs Part 8). Evaluation
followed a \textbf{pre-registered, leakage-controlled OOD protocol}
(threshold calibrated on validation/benign only; OOD = fresh HackAPrompt
injections + dolly benign, deduplicated within and across splits;
touched once). The result required three pre-registered runs, and we
report the full arc --- including two NULLs we held rather than
rationalize --- as the honest record.

\textbf{TABLE VI. Phase-5 head on PG2 frozen embeddings (fresh OOD,
Wilson 95\% CI).}

{\def\LTcaptype{none} % do not increment counter
\begin{table*}[!t]\centering\scriptsize
\resizebox{\textwidth}{!}{%
\begin{tabular}{@{}lllll@{}}
\toprule
Run & OOD Recall & OOD FPR & AUC & Verdict (bar: recall $\ge$50 \%, FPR $\le$1
\%) \\
\midrule



PG2 native head (baseline) & 22.8 \% & 0.25 \% & --- & --- \\
5a (threshold on 45 mismatched benign) & 99.7 \% & 2.2 \% & 1.000 &
\textbf{NULL} (FPR \textgreater{} 1 \%) \\
5a-bis (calibration fixed) & 99.9 \% & 1.2 \% {[}0.8, 1.7{]} & 1.000 &
\textbf{NULL} (point FPR \textgreater{} 1 \%) \\
\textbf{5a-ter (99.5th-pct calib, fresh data)} & \textbf{99.9 \%}
{[}99.3, 100{]} & \textbf{0.7 \%} {[}0.4, 1.2{]} & 0.999 &
\textbf{SUCCESS} \\
\bottomrule
\end{tabular}}
\end{table*}
}

The finding is decisive: a \emph{linear} head on PG2\textquotesingle s
frozen embeddings lifts out-of-distribution injection recall from
\textbf{22.8 \% to 99.9 \% at 0.7 \% FPR}. Phase 4\textquotesingle s
"structural gap" was therefore \textbf{not} in the encoder --- the
representation carries near-perfect signal (AUC $\approx$ 1.0) --- but entirely
in Meta\textquotesingle s precision-first \emph{decision head}. 5a and
5a-bis returned NULL on the hard FPR ceiling (2.2 \%, then 1.2 \% with a
CI that included 1.0 \%); we held both rather than loosen a
pre-registered criterion after seeing the data, and 5a-ter cleared all
three locked bars strictly on fresh disjoint data via a single
pre-registered operating-point change. The composed \textbf{PG2-encoder
$\to$ logistic-head} classifier was then deployed as the origin Tier-3a
model (§IX), replacing the native head. Full protocol and
pre-registration:
\passthrough{\lstinline!docs/phase5a\_finetune\_scoping.md!}.

\subsection{VII-E. Statistical Rigor}\label{vii-e-statistical-rigor}

\begin{itemize}
\tightlist
\item
  \textbf{Sample size for the recall claim.} To credit
  \(\text{Recall}>0.99\) via the Wilson lower bound at 95\% confidence,
  the required \(|\mathcal{X}^{+}|\) follows from the Wilson interval
  width at the observed rate; we state the attained
  \(|\mathcal{X}^{+}|\) and whether it suffices, and if the primary
  corpus is too small for a 0.99 lower bound we say so explicitly rather
  than over-claim from a small sample.
\item
  \textbf{Latency percentile uncertainty.} Tail percentiles from finite
  samples have non-trivial variance; we report bootstrap confidence
  intervals for \(\langle p99\rangle\) and run \(R\ge\) (stated)
  repetitions across time-of-day to bound diurnal network variance.
\item
  \textbf{Multiple comparisons.} When comparing operating points or
  datasets, we note the number of comparisons and avoid selecting the
  best-looking \(\theta\) post hoc without a held-out split.
\end{itemize}

\subsection{VII-F. Threats to Validity}\label{vii-f-threats-to-validity}

\begin{itemize}
\tightlist
\item
  \textbf{Construct / corpus bias.}
  \passthrough{\lstinline!deepset/prompt-injections!} is bilingual DE/EN
  and modestly sized; recall on it may not transfer to other languages
  or to adaptive adversaries. We mitigate by supplementing families
  (§VII-A) and reporting per-family, but do not claim generality beyond
  the tested distribution.
\item
  \textbf{Adaptive adversary.} The evaluation is against a \emph{static}
  corpus; a signature filter is, by construction, evadable by an
  attacker who knows the signatures (novel phrasings).
  §V-D\textquotesingle s scoped claim and the origin-tier escalation
  (§IX) are the response; we do not present static-corpus recall as
  robustness against adaptive attack.
\item
  \textbf{Single-region measurement.} Latency measured from limited
  client geographies may not represent global tail behavior; we report
  the measurement geography and treat \(L_{\text{net}}^{\Delta}\) as
  region-dependent.
\item
  \textbf{Mock origin.} Subtracting a controlled \(S_O\) isolates
  firewall latency but omits real-origin variance interactions; we note
  this and report both firewall-added and end-to-end figures.
\end{itemize}

\subsection{VII-G. Methodological Integrity (the discipline behind the
numbers)}\label{vii-g-methodological-integrity-the-discipline-behind-the-numbers}

The Phase-5 result (Table VI) is a 99.9 \% recall claim, and such claims
are exactly where evaluation most often deceives itself. We therefore
adopted a set of adversarial-to-ourselves controls, and we regard the
\emph{process} as a primary contribution --- a headline number is only
as credible as the protocol that could have falsified it.

\begin{itemize}
\tightlist
\item
  \textbf{Pre-registration.} The full success criterion --- OOD recall
  \(\ge 50\%\) at \emph{point-estimate} FPR \(\le 1\%\), with the recall
  Wilson lower bound clearing the baseline\textquotesingle s upper bound
  --- was fixed in writing
  (\passthrough{\lstinline!docs/phase5a\_finetune\_scoping.md!})
  \textbf{before any training code was written}, and the single
  pre-registered operating-point change for 5a-ter (99.0th \(\to\)
  99.5th calibration percentile) was recorded \textbf{before} the fresh
  draw. The criterion was never renegotiated after seeing data.
\item
  \textbf{Leakage control.} Splits were deduplicated by exact
  normalized-text hash \emph{and} by embedding cosine similarity
  \(\ge 0.95\) \textbf{across} train/validation/OOD (and \(\ge 0.98\)
  within splits), so a reported "win" cannot arise from near-duplicate
  memorization. This removed, e.g., 14 validation items that were
  near-copies of training data, and 137 of 500 sampled OOD injections as
  exact/cross-train matches.
\item
  \textbf{License refusal.} Datasets whose licenses could not be
  confirmed to a primary source were \textbf{excluded from training and
  evaluation} regardless of their utility --- Tensor Trust (license
  unconfirmed; the code repo\textquotesingle s BSD-2 does not cover the
  crowdsourced data), the safe-guard set (no license listed), and BIPIA
  (non-standard license). We accepted a smaller corpus over an
  unverifiable one.
\item
  \textbf{Holding the NULL --- twice.} The decisive discipline: run 5a
  returned 99.7 \% recall at 2.2 \% FPR and run 5a-bis returned 99.9 \%
  recall at 1.2 \% FPR \textbf{with a Wilson CI that included 1.0 \%}. A
  single sentence ("statistically consistent with \(\le 1\%\)") would
  have converted either into a reported SUCCESS. We returned
  \textbf{NULL} both times, because the pre-registered ceiling was on
  the point estimate and loosening it post hoc is precisely the
  self-deception pre-registration exists to prevent. Only after a
  legitimate operating-point change, tested \textbf{once} on genuinely
  fresh, disjoint, never-scored data, did the point estimate clear the
  bar (0.7 \% FPR) --- and \emph{that} number we report as SUCCESS. The
  two NULLs are not failures to hide; they are the evidence that the
  final number was earned, not curated.
\end{itemize}

\begin{center}\rule{0.5\linewidth}{0.5pt}\end{center}

\subsubsection{Citation keys}\label{citation-keys-4}

\begin{itemize}
\tightlist
\item
  \textbf{{[}DS-PI{]}} deepset, "prompt-injections," Hugging Face
  Datasets, huggingface.co/datasets/deepset/prompt-injections (662 rows;
  546/116; Apache 2.0; accessed 2026-07-04).
\item
  \textbf{{[}DS-JB{]}} jackhhao, "jailbreak-classification," Hugging
  Face Datasets (1,306 rows; Apache 2.0; accessed 2026-07-04).
\item
  \textbf{{[}DS-GA{]}} Lakera, "gandalf\_ignore\_instructions," Hugging
  Face Datasets; see also arXiv:2311.01011.
\end{itemize}

\begin{quote}
by direct fetch of the HF dataset cards (verified-facts Part 4b,
2026-07-04). \textbf{Detection results (Table III: edge Recall 2.28 \%,
FPR 0.00 \%, AUC 0.511; Table IV: protectai 41.44 \% / FPR 1.00 \%,
Prompt Guard 2 86M 22.81 \% / FPR 0.25 \%) are MEASURED} over the full
662-row deepset corpus
(\passthrough{\lstinline!benchmarks/results/deepset-662-\{local,fullsystem,promptguard2\}/!})
--- Table III via the firewall\textquotesingle s own
\passthrough{\lstinline!screen()!} pipeline; Table IV\textquotesingle s
protectai row via \passthrough{\lstinline!full\_system.py!}, and its
Prompt Guard 2 row via \passthrough{\lstinline!run\_prompt\_guard.py!}
executed on the \textbf{live production deployment} (Helsinki origin,
CPU). The escalation-tier model specs are CONFIRMED (verified-facts Part
7) but their accuracy figures are vendor-self-reported; Table IV is our
own independent measurement of a proxy. Latency quantities
(\(\langle p50/p90/p99\rangle\), \(\langle\kappa/\rho\rangle\)) remain
placeholders pending an edge deployment. The statistical methods (Wilson
intervals, bootstrap percentiles) are standard; the analytical latency
decomposition (§VII-B) is original and rests only on the CONFIRMED §III
envelope.
\end{quote}

\section{VIII. Related Work}\label{viii-related-work}

We situate the contribution across five strands. Throughout, we cite
only sources whose claims were primary-verified in this work or which
are canonical, dated references; any citation requiring final
bibliographic completion is marked \passthrough{\lstinline![cite]!}.

\textbf{Edge and serverless security.} Web Application Firewalls and
CDN-layer filtering are established, but classical WAF rule engines
target SQL-injection/XSS signatures over HTTP structure, not
natural-language instruction injection into an LLM. The V8-isolate
execution model {[}C-limits{]} gives a materially different cost
envelope than container-per-request FaaS (cold start, per-request
process), which prior serverless-security work assumes; our cost model
(§III-C) is specific to the isolate\textquotesingle s
amortized-construction, CPU-vs-I/O-metered regime and, to our knowledge,
is the first to bound an LLM-injection screener against the
\emph{corrected} (post-2024) Workers CPU envelope rather than the
obsolete 50 ms figure.

\textbf{LLM prompt-injection defenses.} The field spans (i) fine-tuned
structural defenses --- StruQ, SecAlign --- which achieve high
attack-rejection but were recently shown to learn brittle surface
heuristics with large OOD degradation {[}arXiv-2601.07185{]}; (ii)
embedding-based classifiers (RF/ XGBoost over prompt embeddings) that
can outperform encoder-only baselines but at modest absolute AUC on
curated sets {[}arXiv-2410.22284{]}; (iii) model-based guardrails (Llama
Guard-class) \passthrough{\lstinline![cite]!}; and (iv)
deterministic/lexical filters. The OWASP LLM Top 10 {[}OWASP-LLM01{]}
fixes the direct/indirect taxonomy we adopt (§II-C), corroborated by
{[}arXiv-2410.21146{]}. Our position (§V-A) is \emph{not} to compete
with model guardrails on recall but to argue their \emph{architectural}
misplacement at the edge and to occupy the high-precision,
sub-millisecond, I/O-free niche a finite-state matcher uniquely fills
--- with the model tier relocated to the origin (§IX). The multi-pattern
primitive is Aho--Corasick {[}AC75{]}, whose \(O(n+z)\) single-pass
matching {[}Springer-AC{]} is what makes hundreds of signatures
affordable in the isolate budget.

\textbf{Webhook authentication and MAC security.} Production webhook
signing (Stripe\textquotesingle s timestamped
\passthrough{\lstinline!t!}/\passthrough{\lstinline!v1!} HMAC-SHA256
scheme {[}Stripe-sig{]}, GitHub\textquotesingle s equivalent
\passthrough{\lstinline![cite]!}) establishes the practice we formalize.
HMAC itself is RFC 2104 {[}RFC2104{]} / FIPS 198-1 {[}FIPS198-1{]}, with
PRF-security from Bellare--Canetti--Krawczyk {[}BCK96{]} and \emph{tight
multi-user} security from Bellare--Bernstein--Tessaro {[}BBT16{]} ---
the latter is what lets our tenant-binding bound (§IV-C/D) avoid linear
degradation in the tenant count. The novelty is not HMAC but the
\emph{tenant-binding theorem}: to our knowledge, prior webhook-signing
treatments authenticate \emph{message integrity} but do not prove
\emph{tenant-context binding} as a reduction discharging a
cross-tenant-isolation obligation (§II-D $\to$ §IV).

\textbf{Multi-tenant isolation.} Cross-tenant leakage in shared-context
LLM systems is an emerging concern; the adversary model combining a
Dolev--Yao network attacker {[}DY83{]} with a malicious authenticated
tenant (§II-A) is, we believe, the appropriate formalization for
multi-tenant webhook ingestion and is not standard in prior LLM-security
threat models.

\textbf{High-velocity logging.} Time-series and append-only logging over
PostgreSQL is well-trodden operationally; our contribution is the
specific composition --- \passthrough{\lstinline!waitUntil!}-decoupled
writes, a durable queue for at-least-once delivery, JSONB+TOAST for
unstructured payloads, GIN \passthrough{\lstinline!fastupdate!} to bound
index cost, and declarative range partitioning for O(1) retention
{[}PG-toast, PG-gin, PG-partition{]} --- tied to PCI-DSS Req 10.5.1 /
SOC 2 CC7 with edge-side PAN redaction (§VI-F).

\textbf{In sum}, each primitive is prior art; the \emph{contribution is
the vertically-integrated composition} --- edge relocation + a proven
tenant-binding cryptographic layer + a cost-bounded deterministic
screener + a non-blocking compliant logging path --- engineered against
the real, current constraints of the edge runtime and evaluated with
methodological honesty.

\section{IX. Limitations and Future
Work}\label{ix-limitations-and-future-work}

\textbf{Limitations (stated plainly).} (L1) The edge screener is
deterministic and therefore evadable by an adaptive adversary who crafts
novel phrasings outside the signature set; §V-D\textquotesingle s recall
claim is scoped to the signature-expressible class, and static-corpus
recall (§VII) is \emph{not} robustness against adaptive attack. (L2)
HMAC is symmetric, so the scheme provides no \textbf{forward secrecy}: a
leaked per-epoch key forges within its rollover window (§IV-D). (L3)
Cross-tenant \emph{context} isolation at the origin (the
\(c_b\cap d_a=\emptyset\) property of §II-D, mode 1) is \emph{assumed},
not enforced by the edge; the firewall discharges the \emph{attribution}
obligation (mode 2) only. (L4) The replay nonce cache (§IV-E) is
regional; a globally distributed producer set needs distributed replay
state. (L5) Evaluation is single-region, bilingual, and on a
modestly-sized corpus (§VII-F).

\textbf{Future work.}

\begin{enumerate}
\def\labelenumi{\arabic{enumi}.}
\tightlist
\item
  \textbf{Origin-side escalation tier (implemented, deployed,
  measured).} We built and \emph{deployed} the Layer-3 escalation tier
  as a self-hosted classifier and measured it live on the full corpus
  (§VII, Table IV): the chosen multilingual \textbf{Prompt Guard 2 86M}
  achieves 22.8 \% recall at 0.25 \% FPR, and the English-only proxy
  41.4 \% at 1.0 \% FPR --- both an order of magnitude above the edge
  alone, and both far below their vendor self-reports (97.5 \% / 99.7
  \%). We \emph{tested} the two obvious remedies (§VII, Phase 4): a PG2
  threshold sweep lifts recall only to 26.6 \%, and an OR-ensemble with
  NVIDIA NemoGuard (jailbreak axis) adds 0 \% on this injection corpus
  --- compounded by our inability to reproduce
  NemoGuard\textquotesingle s CPU embedding pipeline from public
  artifacts (a blatant jailbreak scored \textasciitilde0.09; we do not
  claim NemoGuard is weak, only that we could not faithfully run it
  off-NIM). That the flagship model sits at 22.8 \% recall and neither
  of those two remedies closes the gap was, at Phase 4, the key open
  question. \textbf{Phase 5 resolved it (§VII, Table VI): the gap is NOT
  structural.} A linear head trained on PG2\textquotesingle s
  \emph{frozen} embeddings reaches \textbf{99.9 \% OOD recall at 0.7 \%
  FPR} (AUC $\approx$ 1.0) under a pre-registered, leakage-controlled protocol
  --- proving the encoder always carried near-perfect signal and the
  22.8 \% was purely Meta\textquotesingle s precision-first decision
  head. The composed \textbf{PG2-encoder $\to$ logistic-head} classifier is
  deployed as the origin Tier-3a model, replacing the native head.
  Remaining future work: a cost model bounding the escalation rate,
  adaptive-attack hardening, and periodic re-calibration of the head
  threshold on deployment traffic.
\item
  \textbf{Threat-log-driven signature synthesis (closed loop).} Mine the
  §VI JSONB threat log to auto-propose new Aho--Corasick signatures from
  clustered novel payloads, with human-in-the-loop validation ---
  turning the logging tier from a passive record into an active
  detection feedback loop, and partially closing L1 against
  slowly-adapting adversaries.
\item
  \textbf{Forward-secure / asymmetric webhook signing.} Replace or
  augment symmetric HMAC with per-epoch key ratcheting or asymmetric
  signatures to obtain forward secrecy and reduce key-distribution
  burden (addressing L2), quantifying the added edge CPU cost of
  asymmetric verification against the §III budget.
\item
  \textbf{TEE-backed origin context isolation.} Enforce (rather than
  assume) the mode-1 isolation of L3 with hardware-isolated per-tenant
  model contexts, extending the cryptographic guarantee across trust
  boundary B3.
\item
  \textbf{Distributed replay state.} Realize the nonce cache as
  globally-consistent state (e.g., Durable Objects) with an explicit
  latency/consistency trade characterization (addressing L4).
\item
  \textbf{Adaptive-adversary and multi-modal evaluation.} Extend the
  corpus to adaptive/obfuscated and non-text (image/file) injection
  vectors, and characterize the economic ("denial-of-wallet") DoS bound
  formally.
\end{enumerate}

\section{X. Conclusion}\label{x-conclusion}

We presented a multi-layered defensive architecture that relocates the
first line of LLM-webhook defense to the network edge and integrates
four contributions into one request path: (i) an edge-isolate firewall
whose per-request cost is provably bounded --- strictly linear in a
capped input, backtrack-free, with one-time automaton construction
amortized to module scope --- and analyzed against the \emph{corrected}
Cloudflare Workers envelope (10 ms free / 30 s--5 min paid CPU, 128 MB
isolate), replacing the widely-repeated but obsolete "50 ms" premise
with primary-sourced facts; (ii) a cryptographic tenant-isolation layer
whose tenant-binding property is proven by reduction to
HMAC\textquotesingle s multi-user EUF-CMA security, yielding a bound
independent of tenant count and key-rotation multiplicity, with replay
handled as an orthogonal freshness property; (iii) a deterministic
Aho--Corasick screening layer, honestly positioned as a high-precision
sub-millisecond first filter whose recall gap on novel semantic attacks
is delegated to an origin-side model tier; and (iv) a non-blocking,
compliance-aware threat-logging path over partitioned PostgreSQL JSONB
that never gates the request and never persists cardholder data.

Two commitments distinguish the work methodologically: every platform
figure is primary-sourced and date-stamped against an actively-changing
edge platform, and every \emph{result} quantity is deferred to a
reproducible benchmark harness rather than asserted --- the paper
specifies the experiment and lets measurement dictate the claims.

\textbf{The deployed system, operating as the AIO Apex enterprise
firewall, is a dual-tier composition.} A multi-tenant Cloudflare edge
performs the rapid, sub-millisecond screen --- HMAC-SHA256 tenant
binding for authenticated isolation, plus Aho--Corasick signature
matching --- and \textbf{fails open to a self-hosted origin
deep-scanner} when the model tier is unreachable, so availability never
depends on the heavier layer. That origin tier now runs the Phase-5
classifier: \textbf{Prompt Guard 2\textquotesingle s frozen encoder
feeding a calibrated linear head}.

\textbf{The central empirical result is that the recall gap was never
structural.} Meta\textquotesingle s Prompt Guard 2 exhibits 22.8 \%
out-of-distribution recall through its \emph{native} head, and Phase 4
might have concluded the encoder was blind to those attacks. Phase 5
disproves that: a \emph{linear} logistic head on the
model\textquotesingle s frozen penultimate embeddings reaches
\textbf{99.9 \% OOD recall at 0.7 \% false-positive rate (AUC
\(\approx 0.999\))} under a pre-registered, leakage-controlled protocol
on fresh, disjoint data (Table VI). The 22.8 \% was purely a
\textbf{decision-head artifact} --- Meta\textquotesingle s
precision-first calibration --- not an encoder limitation; the injection
signal was present in the representation all along and merely required a
properly calibrated boundary to extract. This composed head is deployed
as the origin tier, replacing the native head end-to-end.

Equally, we regard the \emph{discipline} that produced that number as a
contribution in its own right (§VII-G): success criteria pre-registered
before any code, cosine-\(\ge 0.95\) cross-split deduplication, refusal
of unverifiable-license data, and --- decisively --- \textbf{two NULL
verdicts held on 99.7 \%+ recall} because a strict FPR point-estimate
ceiling was not cleared, before the bar was met on genuinely fresh data.
A 99.9 \% that survives that process is worth more than one asserted
without it. The architecture demonstrates that strong, provable
multi-tenant LLM webhook security --- and a detector that closes the
recall gap --- is achievable within the real constraints of edge
compute, with vast CPU headroom rather than at its margin, provided each
layer is engineered to the platform\textquotesingle s actual, current
envelope and each claim is earned rather than curated.

\begin{center}\rule{0.5\linewidth}{0.5pt}\end{center}

\subsubsection{Additional citation keys (canonical /
dated)}\label{additional-citation-keys-canonical--dated}

\begin{itemize}
\tightlist
\item
  \textbf{{[}AC75{]}} A. V. Aho, M. J. Corasick, "Efficient String
  Matching: An Aid to Bibliographic Search," Communications of the ACM,
  18(6):333--340, 1975.
\item
  \textbf{{[}DY83{]}} D. Dolev, A. C. Yao, "On the Security of Public
  Key Protocols," IEEE Trans. Information Theory, 29(2):198--208, 1983.
\item
  \textbf{{[}Wilson27{]}} E. B. Wilson, "Probable Inference, the Law of
  Succession, and Statistical Inference," J. American Statistical
  Association, 22(158):209--212, 1927. \emph{(§VII-E interval.)}
\item
  (Prior citation keys {[}C-limits{]}, {[}OWASP-LLM01{]},
  {[}arXiv-2601.07185{]}, {[}arXiv-2410.22284{]},
  {[}arXiv-2410.21146{]}, {[}Springer-AC{]}, {[}Stripe-sig{]},
  {[}RFC2104{]}, {[}FIPS198-1{]}, {[}BCK96{]}, {[}BBT16{]},
  {[}PG-\emph{{]}, {[}DS-}{]} resolve as defined in their home
  sections.)
\end{itemize}

\begin{quote}
1975, Dolev--Yao 1983, Wilson 1927 --- all real, well-established);
items needing final bibliographic detail (Llama Guard, GitHub webhook
docs) are marked \passthrough{\lstinline![cite]!} and must be completed
before submission --- none is load-bearing for a claim. §IX limitations
are the honest consolidation of the non-claims made throughout (forward
secrecy, adaptive adversary, assumed origin isolation). §X asserts no
new facts.
\end{quote}

\section{Appendix A. Full STRIDE $\to$ Mitigation
Mapping}\label{appendix-a-full-stride--mitigation-mapping}

Table A.1 expands the abbreviated Table I (§II-B) into the complete
decomposition, adding the cross-boundary coupling, a qualitative
likelihood/impact for the payment-gateway running example, the
mitigating control with its paper section, and the residual risk that
survives the control. "L/I" is Likelihood/Impact on a \{Low, Med, High\}
scale. Boundaries: \textbf{B1} public network, \textbf{B2} edgeorigin,
\textbf{B3} the model\textquotesingle s instruction/data boundary
(§II-A).

\textbf{TABLE A.1. Complete STRIDE decomposition of the webhook$\to$LLM
pipeline.}

{\def\LTcaptype{none} % do not increment counter
\begin{table*}[!t]\centering\scriptsize
\resizebox{\textwidth}{!}{%
\begin{tabular}{@{}lllllllll@{}}
\toprule
\# & STRIDE & Threat & Asset & Bdy & Coupling & L/I & Mitigating control
(§) & Residual risk \\
\midrule



1 & Spoofing & Forged/unsigned webhook impersonating a producer or
tenant & A1 & B1 & --- & H/H & HMAC-SHA256 verify + tenant binding
(§IV-C) & Key compromise within a rollover window (§IV-D) \\
2 & Tampering & Payload mutated in transit & A3 & B1 & --- & M/H &
Signature over canonical \passthrough{\lstinline!H\_body!} (§IV-B) &
None if HMAC unbroken \\
3 & Tampering$\to$Elev. & Indirect injection: adversarial instructions in
relayed content & A3 & B1,B3 & T$\to$E & H/H & Aho--Corasick + feature
screening (§V) & Novel/semantic phrasings (measured FNR 97.7 \%, §VII);
delegated to origin tier (§IX) \\
4 & Tampering$\to$Info & Cross-tenant context poisoning: \(t_a\) mutates
shared state read into \(t_b\)\textquotesingle s context & A2,A3 & B3 &
T$\to$I & L/H & Tenant binding (§IV) + assumed per-tenant context isolation
(§II-D mode 1) & Origin-side isolation is \emph{assumed}, not enforced
by the edge (L3, §IX) \\
5 & Repudiation & Attack occurs with no durable attribution & A5 & B1 &
--- & M/M & Async append-only threat log (§VI); PCI 10.2/10.3 &
Log-store availability over the write window (§VI-B) \\
6 & Info. disclosure & System-prompt / context extraction; cross-tenant
leakage & A2 & B3 & --- & M/H & Extraction signatures (§V-B) + isolation
& Obfuscated extraction below signature threshold \\
7 & DoS & Token flood / CPU / subrequest exhaustion; "denial-of-wallet"
& A4 & B1 & --- & H/M & Size guard + sub-ms rejection before inspection
(§III-C, §V) & Distributed volumetric floods (upstream CDN scope) \\
8 & Elev. of priv. & Direct injection / jailbreak executes with the
agent\textquotesingle s tool/data privileges & A2,A3 & B3 & --- & H/H &
Screening (§V) + least-privilege origin tools & Same recall gap as row
3 \\
\bottomrule
\end{tabular}}
\end{table*}
}

\textbf{Reading the residual-risk column.} Rows 3 and 8 carry the recall
gap the paper measures empirically (§VII, Table III): the edge
layer\textquotesingle s deterministic screening is high-precision but
low-recall, so the \emph{residual} for these rows is explicitly
transferred to the origin escalation tier (§IX), not claimed as closed.
Row 4\textquotesingle s residual --- that mode-1 context isolation is
assumed --- is the paper\textquotesingle s principal scoping boundary
(L3). This column is what makes the threat model \emph{actionable}: it
states, per threat, exactly what is and is not discharged.

\begin{center}\rule{0.5\linewidth}{0.5pt}\end{center}

\section{Appendix B. Proof of the Tenant-Binding
Theorem}\label{appendix-b-proof-of-the-tenant-binding-theorem}

We give the full game-hopping proof of Theorem 1 (§IV-C), the
Injectivity Lemma it depends on, and the extensions to key rotation
(Theorem 2) and replay.

\subsection{B.1 Preliminaries recalled}\label{b1-preliminaries-recalled}

Let \(H=\mathsf{HMAC\text{-}SHA256}\) with tag length \(n=256\). For a
set of tenants \(\mathcal{T}\), keys \(k_i \leftarrow_\$ \{0,1\}^n\) are
drawn independently. The signed message is the injective encoding
\fiteq{m = \langle \mathsf{tid} \rangle \,\|\, \langle \kappa \rangle \,\|\, \langle t \rangle
\,\|\, \langle \eta \rangle \,\|\, \langle H_{\text{body}} \rangle,}
where \(\langle\cdot\rangle\) denotes a length-prefixed
(self-delimiting) field encoding. Write \(\mathsf{tid}(m)\) for the
tenant-id field recovered from \(m\). The signing function is
\(\mathsf{Mac}_{k}(m)=H_k(m)\) and verification checks
\(H_k(m)\stackrel{?}{=}\tau\) in constant time.

\textbf{Assumptions.} (i) Multi-user PRF security of HMAC: for \(u\)
instances, \(\mathbf{Adv}^{\text{mu-prf}}_{H,u}(\mathcal{B})\) is
negligible and, by {[}BBT16{]}, carries no factor linear in \(u\). (ii)
The field encoding is injective (Lemma B.1).

\subsection{\texorpdfstring{B.2 The multi-user tenant-binding game
\(\mathbf{G}_0\)}{B.2 The multi-user tenant-binding game \textbackslash mathbf\{G\}\_0}}\label{b2-the-multi-user-tenant-binding-game-mathbfg_0}

\begin{enumerate}
\def\labelenumi{\arabic{enumi}.}
\tightlist
\item
  Sample \(k_i\leftarrow_\$\{0,1\}^n\) for all \(t_i\in\mathcal{T}\).
\item
  \(\mathcal{A}\) chooses a corrupt set
  \(\mathcal{C}\subsetneq\mathcal{T}\) and receives
  \(\{k_i: t_i\in\mathcal{C}\}\). Let
  \(\mathcal{U}=\mathcal{T}\setminus\mathcal{C}\), \(u=|\mathcal{U}|\).
\item
  \(\mathcal{A}\) has oracles, for every \(t_j\in\mathcal{U}\):
  \(\mathsf{Sign}_j(m)=H_{k_j}(m)\) (queries recorded in
  \(\mathcal{Q}_j\)) and
  \(\mathsf{Ver}_j(m,\tau)=[\,H_{k_j}(m)=\tau\,]\).
\item
  \(\mathcal{A}\) outputs \((b,m^*,\tau^*)\) with
  \(t_b\in\mathcal{U}\) and \(\mathsf{tid}(m^*)=\mathsf{tid}_b\).
\item
  \textbf{Win} (\(\mathbf{G}_0=1\)) iff \(H_{k_b}(m^*)=\tau^*\) and
  \(m^*\notin\mathcal{Q}_b\).
\end{enumerate}

By definition
\(\mathbf{Adv}^{\text{bind}}(\mathcal{A})=\Pr[\mathbf{G}_0=1]\). Let
\(Q_v\) be the total number of verification queries \(\mathcal{A}\)
makes.

\subsection{B.3 Lemma B.1 (Injectivity / namespace
separation)}\label{b3-lemma-b1-injectivity--namespace-separation}

\emph{If the field encoding \(\langle\cdot\rangle\|\cdots\) is injective
and \(\mathsf{tid}\) occupies a self-delimiting field, then (a) distinct
field-tuples yield distinct \(m\), and (b) any \(m^*\) with
\(\mathsf{tid}(m^*)=\mathsf{tid}_b\) that is not in \(\mathcal{Q}_b\)
is a fresh input to \(H_{k_b}\) --- independent of every query made to
any \(\mathsf{Sign}_j\), \(j\neq b\).}

\emph{Proof.} (a) is immediate from injectivity of a concatenation of
self-delimiting fields. For (b): the only oracle evaluating \(H_{k_b}\)
is \(\mathsf{Sign}_b\) (and \(\mathsf{Ver}_b\), which reveals only a
bit). Queries to \(\mathsf{Sign}_j\), \(j\neq b\), evaluate \(H_{k_j}\)
under an independent key \(k_j\) and thus reveal nothing about
\(H_{k_b}(m^*)\). Since \(m^*\notin\mathcal{Q}_b\), the value
\(H_{k_b}(m^*)\) was never returned to \(\mathcal{A}\). \(\square\)

The lemma is the crux of \emph{binding}: because \(\mathsf{tid}_b\) is
inside \(m\), owning every other tenant\textquotesingle s key (hence
every \(H_{k_j}\), \(j\neq b\)) yields no information about the tag
\(\mathcal{A}\) must produce for \(t_b\).

\subsection{B.4 Proof of Theorem 1}\label{b4-proof-of-theorem-1}

We bound \(\Pr[\mathbf{G}_0=1]\) by a two-game hop.

\textbf{Game \(\mathbf{G}_1\).} Identical to \(\mathbf{G}_0\), except
each uncorrupted tenant\textquotesingle s \(H_{k_j}\)
(\(j\in\mathcal{U}\)) is replaced by an independent truly random
function \(R_j:\{0,1\}^*\to\{0,1\}^n\); both \(\mathsf{Sign}_j\) and
\(\mathsf{Ver}_j\) use \(R_j\).

\emph{Claim 1.} \(\big|\Pr[\mathbf{G}_1=1]-\Pr[\mathbf{G}_0=1]\big|\le
\mathbf{Adv}^{\text{mu-prf}}_{H,u}(\mathcal{B})\).

\emph{Proof.} Construct a mu-PRF distinguisher \(\mathcal{B}\) against
\(u\) instances. \(\mathcal{B}\) samples the corrupt keys itself,
answers corrupt-key reveals directly, and routes every
\(\mathsf{Sign}_j/\mathsf{Ver}_j\) call (\(j\in\mathcal{U}\)) to its
\(j\)-th oracle. If the oracles are the real \(H_{k_j}\),
\(\mathcal{B}\) simulates \(\mathbf{G}_0\) exactly; if they are random
functions, it simulates \(\mathbf{G}_1\). \(\mathcal{B}\) outputs \(1\)
iff \(\mathcal{A}\) wins. The distinguishing advantage is the game gap,
and by {[}BBT16{]} it is \(\le\mathbf{Adv}^{\text{mu-prf}}_{H,u}\) with
no \(u\)-factor. \(\square\)

\emph{Claim 2.} \(\Pr[\mathbf{G}_1=1]\le Q_v\,2^{-n}\).

\emph{Proof.} In \(\mathbf{G}_1\), verification of a candidate
\((b,m^*,\tau^*)\) succeeds iff \(\tau^*=R_b(m^*)\). By Lemma
B.1(b), \(m^*\) (fresh, \(\mathsf{tid}=\mathsf{tid}_b\),
\(\notin\mathcal{Q}_b\)) is a point at which \(R_b\) has never been
evaluated in \(\mathcal{A}\)\textquotesingle s view, so \(R_b(m^*)\) is
uniform on \(\{0,1\}^n\) and independent of everything \(\mathcal{A}\)
has seen (including all \(R_{j\neq b}\) outputs and all corrupt keys).
Any single verification query therefore succeeds with probability
exactly \(2^{-n}\). A union bound over the \(\le Q_v\) verification
queries gives \(\Pr[\mathbf{G}_1=1]\le Q_v\,2^{-n}\). \(\square\)

Combining,
\(\mathbf{Adv}^{\text{bind}}(\mathcal{A})=\Pr[\mathbf{G}_0=1]\le
\mathbf{Adv}^{\text{mu-prf}}_{H,u}(\mathcal{B})+Q_v\,2^{-n}\), which is
Theorem 1. \(\blacksquare\)

\textbf{Interpretation.} With \(n=256\), the \(Q_v\,2^{-n}\) term is
negligible for any feasible \(Q_v\); the bound is dominated by the
mu-PRF term, which --- critically --- does not grow with the number of
tenants \(u\). Cross-tenant attribution is therefore as hard as breaking
HMAC, at enterprise tenant scale.

\subsection{B.5 Extension to key rotation (Theorem
2)}\label{b5-extension-to-key-rotation-theorem-2}

Model each currently-valid pair \((t_i,\kappa)\), \(t_i\in\mathcal{U}\),
as one instance; let
\(u'=\sum_{t_i\in\mathcal{U}}|\{\kappa: k_{i,\kappa}\text{ valid}\}|\le uL\).
The rotation verifier accepts a \(t_b\)-attributed message iff some
valid \((t_b,\kappa)\) verifies it, i.e. iff it is a forgery under one
of these instances. Since \(\kappa\) is inside \(m\) (authenticated),
Lemma B.1 extends: a message with a given \((\mathsf{tid}_b,\kappa)\) is
fresh input to instance \((b,\kappa)\). Re-running §B.4 with the
instance set enlarged from \(u\) to \(u'\) gives
\(\mathbf{Adv}^{\text{bind-rot}}\le\mathbf{Adv}^{\text{mu-prf}}_{H,u'}+Q_v\,2^{-n}\),
which by {[}BBT16{]} is again independent of \(u'\) (hence of both \(u\)
and \(L\)). \(\blacksquare\)

\subsection{B.6 Replay resistance (orthogonal freshness
argument)}\label{b6-replay-resistance-orthogonal-freshness-argument}

Theorem 1 permits replay of an \emph{already-signed} \((m,\tau)\):
\(m\in\mathcal{Q}_b\) is not a win. Freshness is enforced separately.
With tolerance window \(\Delta\), nonce cache \(\mathcal{N}\) (TTL
\(2\Delta\)), and \(\lambda\)-bit nonces, a replayed \((m,\tau)\) is
accepted only if it arrives within \(\Delta\) of \(t_s\) (else the
timestamp check rejects) \emph{and} its nonce \(\eta\) has been evicted
from \(\mathcal{N}\) (else the replay check rejects). Since
\(\mathcal{N}\) retains \(\eta\) for \(2\Delta\ge\Delta\), within the
timestamp window \(\eta\) is necessarily present, so
\fiteq{\Pr[\text{replay accepted}]\le\Pr[\text{nonce-store loss over }2\Delta]+2^{-\lambda},}
where the first term is the cache\textquotesingle s availability failure
probability and \(2^{-\lambda}\) bounds an accidental nonce collision.
For \(\lambda=128\) the collision term is negligible; replay resistance
reduces to nonce-store reliability. \(\blacksquare\)

\subsection{B.7 Remarks}\label{b7-remarks}

\begin{itemize}
\tightlist
\item
  \textbf{Tightness.} The reduction is tight up to the mu-PRF term;
  there is no target-guessing \(1/u\) loss (a forgery under \emph{any}
  uncorrupted instance wins the mu-game).
\item
  \textbf{No forward secrecy.} Keys are symmetric; a leaked
  \(k_{i,\kappa}\) forges within \(\kappa\)\textquotesingle s window.
  Rotation bounds, but does not retroactively protect, that window.
\item
  \textbf{Constant-time verification} is required so that
  \(\mathsf{Ver}\) leaks only the accept/reject bit modelled in the
  game, not tag bytes via timing.
\end{itemize}

\begin{center}\rule{0.5\linewidth}{0.5pt}\end{center}

\section{Appendix C. Signature Corpus and
Reproducibility}\label{appendix-c-signature-corpus-and-reproducibility}

\subsection{C.1 Signature families}\label{c1-signature-families}

The shipped corpus
(\passthrough{\lstinline!edge-firewall/src/signatures.ts!}) contains 24
signatures across the OWASP LLM01 families of §II-C. Each is
\passthrough{\lstinline!blocking!} (a single hit forces reject,
favouring the low-FPR operating point) or weighted (contributes to the
soft score \(\theta\), §V-C).

\textbf{TABLE C.1. Signature families (shipped corpus).}

{\def\LTcaptype{none} % do not increment counter
\begin{table*}[!t]\centering\scriptsize
\resizebox{\textwidth}{!}{%
\begin{tabular}{@{}llll@{}}
\toprule
Family (§II-C) & Example signatures & Blocking & Weighted \\
\midrule



Direct override &
\passthrough{\lstinline!ignore previous instructions!},
\passthrough{\lstinline!disregard all previous!},
\passthrough{\lstinline!forget everything above!} & 3 & 2 \\
Extraction & \passthrough{\lstinline!reveal your system prompt!},
\passthrough{\lstinline!repeat the words above!},
\passthrough{\lstinline!print your configuration!} & 1 & 3 \\
Role-play / persona & \passthrough{\lstinline!you are now dan!},
\passthrough{\lstinline!enable developer mode!},
\passthrough{\lstinline!do anything now!} & 2 & 3 \\
Safety bypass & \passthrough{\lstinline!jailbreak mode!},
\passthrough{\lstinline!without any content filter!},
\passthrough{\lstinline!unlock full capabilities!} & 1 & 2 \\
Delimiter / tag & `\textless{} & system &
\textgreater{}\passthrough{\lstinline!, !}\textless{} \\
Control-token & `\textless{} & endoftext &
\textgreater{}\passthrough{\lstinline!, !}` \\
\bottomrule
\end{tabular}}
\end{table*}
}

The compactness of this set is deliberate and is exactly what §VII,
Table III measures: it yields 0 \% FPR but low recall on diverse real
injections. The set is \emph{not} tuned against the evaluation corpus
(that would overfit and is methodologically disallowed); expanding it is
future work via threat-log-driven synthesis (§IX-2).

\subsection{C.2 Evaluation corpus}\label{c2-evaluation-corpus}

\begin{itemize}
\tightlist
\item
  \textbf{Source:} \passthrough{\lstinline!deepset/prompt-injections!},
  Hugging Face, Apache 2.0.
\item
  \textbf{Size:} 662 rows (546 train / 116 test); labels \(0=\)legit,
  \(1=\)injection; DE+EN.
\item
  \textbf{Split used:} the full 662 rows (\(|\mathcal{X}^+|=263\),
  \(|\mathcal{X}^-|=399\)).
\item
  \textbf{Fetched by:}
  \passthrough{\lstinline!benchmarks/corpus/fetch\_corpus.mjs!}
  (datasets-server rows API).
\end{itemize}

\subsection{C.3 Operating point and
metrics}\label{c3-operating-point-and-metrics}

Decision threshold \(\theta=1.0\) (§V-C). A payload is a predicted
positive iff the firewall does not \passthrough{\lstinline!forward!}
(i.e. \passthrough{\lstinline!reject!} or
\passthrough{\lstinline!escalate!}). Metrics as defined in §VII-D;
confidence intervals are Wilson score at 95 \%
(\passthrough{\lstinline!benchmarks/lib/stats.mjs!}); ROC AUC by
trapezoidal integration over the swept threshold, with a blocking hit
treated as a maximal-confidence positive.

\subsection{C.4 Reproduction}\label{c4-reproduction}

\begin{lstlisting}[language=bash]
cd benchmarks && npm install
npm run corpus      # deepset/prompt-injections -> corpus/payloads.json (662 rows)
npm run detection   # confusion matrix + Wilson CIs  -> results/detection.{json,md}
npm run roc         # ROC/AUC over the real screen()  -> results/roc.json
npm run report      # Section VII placeholder table
\end{lstlisting}

Committed artifacts of the run reported in Table III live in
\passthrough{\lstinline!benchmarks/results/deepset-662-local/!}. Latency
(§VII-C) requires a live edge deployment; run
\passthrough{\lstinline!benchmarks/latency.k6.js!} with
\passthrough{\lstinline!WORKER\_URL!} and a provisioned bench tenant
key.

\begin{quote}
residuals. Appendix B reduces solely to the mu-PRF assumption
({[}BBT16{]}) and the standard EUF-CMA framework --- no unverified
external facts. Appendix C\textquotesingle s corpus facts are CONFIRMED
by direct fetch (verified-facts Part 4b) and its metrics are the
executed harness output.
\end{quote}

\begin{center}\rule{0.5\linewidth}{0.5pt}\end{center}

\section{Data and Code Availability}\label{data-and-code-availability}

The Phase-5 injection-classifier pipeline is released to support
reproduction of the 99.9 \% out-of-distribution result: the corpus
builder and leakage-control deduplication
(\passthrough{\lstinline!build\_corpus.py!}), the frozen Prompt Guard 2
embedding extractor
(\passthrough{\lstinline!extract\_pg2\_embeddings.py!}), the
head-training and single-touch OOD evaluation
(\passthrough{\lstinline!train\_head.py!},
\passthrough{\lstinline!phase5a\_bis.py!},
\passthrough{\lstinline!phase5a\_ter.py!}), the production head
finalizer (\passthrough{\lstinline!finalize\_head.py!}), and the
pre-registration protocol
(\passthrough{\lstinline!docs/phase5a\_finetune\_scoping.md!}) that
fixed the success criteria before any training code. All evaluation
datasets are the license-confirmed public corpora enumerated in Appendix
C. The deterministic seeds (20260705/06/07) reproduce the exact train /
validation / OOD splits.

Repository:
\textbackslash url\{\url{https://github.com/mosafariuk/prompt-guard-2-frozen-head\%7D}
\% AIO Apex --- set final URL

Artifacts are additionally available from the corresponding author
(\href{mailto:mo@aioapex.com}{\nolinkurl{mo@aioapex.com}}) on request.

\section{References}\label{references}

Consolidated bibliography. Keys are the mnemonic
\passthrough{\lstinline!\\cite!} labels used inline throughout the
sections; a camera-ready build would renumber these {[}1{]}\ldots{[}n{]}
via BibTeX. Grouped by type for readability; \textbf{P} =
primary/verified, \textbf{C} = canonical dated work, \textbf{S} =
secondary (used only qualitatively).

\subsection{Platform documentation (P --- Cloudflare, accessed
2026-07-04)}\label{platform-documentation-p--cloudflare-accessed-2026-07-04}

\begin{itemize}
\tightlist
\item
  \textbf{{[}C-limits{]}} Cloudflare, "Workers --- Limits,"
  developers.cloudflare.com/workers/platform/limits.
\item
  \textbf{{[}C-pricing{]}} Cloudflare, "Workers --- Pricing,"
  developers.cloudflare.com/workers/platform/pricing.
\item
  \textbf{{[}C-changelog-2025{]}} Cloudflare, "Higher CPU limits for
  Workers," changelog, 2025-03-25.
\item
  \textbf{{[}C-changelog-2026{]}} Cloudflare, "Increased subrequest
  limits," changelog, 2026-02-11.
\item
  \textbf{{[}C-do-limits{]}} Cloudflare, "Durable Objects --- Limits,"
  developers.cloudflare.com/durable-objects/platform/limits.
\item
  \textbf{{[}C-wf-limits{]}} Cloudflare, "Workflows --- Limits,"
  developers.cloudflare.com/workflows/reference/limits.
\end{itemize}

\subsection{Cryptography (P/C)}\label{cryptography-pc}

\begin{itemize}
\tightlist
\item
  \textbf{{[}RFC2104{]}} H. Krawczyk, M. Bellare, R. Canetti, "HMAC:
  Keyed-Hashing for Message Authentication," RFC 2104, IETF, Feb. 1997.
\item
  \textbf{{[}FIPS198-1{]}} NIST, "The Keyed-Hash Message Authentication
  Code (HMAC)," FIPS PUB 198-1, 2008.
\item
  \textbf{{[}BCK96{]}} M. Bellare, R. Canetti, H. Krawczyk, "Keying Hash
  Functions for Message Authentication," CRYPTO 1996. (See also M.
  Bellare, "New Proofs for NMAC and HMAC," CRYPTO 2006.)
\item
  \textbf{{[}BBT16{]}} M. Bellare, D. J. Bernstein, S. Tessaro,
  "Hash-Function Based PRFs: AMAC and Its Multi-User Security,"
  EUROCRYPT 2016.
\item
  \textbf{{[}Stripe-sig{]}} Stripe, "Verify webhook signatures,"
  docs.stripe.com/webhooks/signature (accessed 2026-07-04).
\end{itemize}

\subsection{LLM security \& prompt injection
(P)}\label{llm-security--prompt-injection-p}

\begin{itemize}
\tightlist
\item
  \textbf{{[}OWASP-LLM01{]}} OWASP, "LLM01:2025 Prompt Injection," Top
  10 for LLM Applications,
  genai.owasp.org/llmrisk/llm01-prompt-injection.
\item
  \textbf{{[}arXiv-2601.07185{]}} "Defenses Against Prompt Attacks Learn
  Surface Heuristics," arXiv:2601.07185, 2026.
\item
  \textbf{{[}arXiv-2410.22284{]}} M. A. Ayub, S. Majumdar,
  "Embedding-based classifiers for prompt-injection detection," CAMLIS
  2024, arXiv:2410.22284.
\item
  \textbf{{[}arXiv-2410.21146{]}} "(direct/indirect/stored
  prompt-injection taxonomy)," arXiv:2410.21146, 2024.
\end{itemize}

\subsection{Algorithms \& statistics
(C)}\label{algorithms--statistics-c}

\begin{itemize}
\tightlist
\item
  \textbf{{[}AC75{]}} A. V. Aho, M. J. Corasick, "Efficient String
  Matching: An Aid to Bibliographic Search," Comm. ACM, 18(6):333--340,
  1975.
\item
  \textbf{{[}Springer-AC{]}} "(Aho--Corasick multi-pattern matching for
  signature detection)," Springer, doi:10.1007/978-3-031-96093-2\_15.
\item
  \textbf{{[}DY83{]}} D. Dolev, A. C. Yao, "On the Security of Public
  Key Protocols," IEEE Trans. Inf. Theory, 29(2):198--208, 1983.
\item
  \textbf{{[}Wilson27{]}} E. B. Wilson, "Probable Inference, the Law of
  Succession, and Statistical Inference," J. Amer. Stat. Assoc.,
  22(158):209--212, 1927.
\end{itemize}

\subsection{Database (P --- PostgreSQL official
docs)}\label{database-p--postgresql-official-docs}

\begin{itemize}
\tightlist
\item
  \textbf{{[}PG-json{]}} "JSON Types,"
  postgresql.org/docs/current/datatype-json.html.
\item
  \textbf{{[}PG-toast{]}} "Database Physical Storage --- TOAST,"
  postgresql.org/docs/current/storage-toast.html.
\item
  \textbf{{[}PG-wal-conf{]}} "Write Ahead Log --- runtime config,"
  postgresql.org/docs/current/runtime-config-wal.html.
\item
  \textbf{{[}PG-wal-config{]}} "WAL Configuration,"
  postgresql.org/docs/current/wal-configuration.html.
\item
  \textbf{{[}PG-gin{]}} "GIN Indexes,"
  postgresql.org/docs/current/gin.html.
\item
  \textbf{{[}PG-partition{]}} "Table Partitioning,"
  postgresql.org/docs/current/ddl-partitioning.html.
\end{itemize}

\subsection{Compliance (P)}\label{compliance-p}

\begin{itemize}
\tightlist
\item
  \textbf{{[}PCI-DSS{]}} PCI Security Standards Council, "PCI-DSS
  v4.0.1," June 2024. Req 10.2/10.3.1/10.3.2/10.5.1; Req
  3.3/3.4.1/3.5.1.
\item
  \textbf{{[}AICPA-TSC{]}} AICPA, "Trust Services Criteria" (2017, rev.
  2022), Common Criteria CC7.1--CC7.4.
\end{itemize}

\subsection{Datasets (P --- accessed
2026-07-04)}\label{datasets-p--accessed-2026-07-04}

\begin{itemize}
\tightlist
\item
  \textbf{{[}DS-PI{]}} deepset, "prompt-injections," Hugging Face
  Datasets (662 rows; 546/116; Apache 2.0).
\item
  \textbf{{[}DS-JB{]}} jackhhao, "jailbreak-classification," Hugging
  Face Datasets (1,306 rows; Apache 2.0).
\item
  \textbf{{[}DS-GA{]}} Lakera, "gandalf\_ignore\_instructions," Hugging
  Face Datasets; see also arXiv:2311.01011.
\end{itemize}

\subsection{Secondary (S --- qualitative use
only)}\label{secondary-s--qualitative-use-only}

\begin{itemize}
\tightlist
\item
  \textbf{{[}Layered{]}} "(layered prompt-injection detector; \textless1
  ms Layers 1--2, 300--800 ms model layer)," dev.to. Used only for the
  model-tier latency order-of-magnitude (§V-A/§V-D).
\end{itemize}

\subsection{To complete before
camera-ready}\label{to-complete-before-camera-ready}

\begin{itemize}
\tightlist
\item
  \textbf{{[}cite{]}} Llama Guard (Meta) --- model-guardrail reference
  (§V-A, §VIII).
\item
  \textbf{{[}cite{]}} GitHub webhook signing docs --- timestamped HMAC
  scheme (§VIII).
\end{itemize}

%% -----------------------------------------------------------------------------

%% Bibliography: convert paper/10_references.md to refs.bib, or use \thebibliography.
% \bibliographystyle{IEEEtran}
% \bibliography{refs}

\end{document}
